\chapter{Objectifs d'apprentissage}\label{IV-objectifs-dapprentissage}

Cette section s'efforce d'expliciter et de hiérarchiser les acquis
d'apprentissages associés au chapitre. Ces objectifs sont organisés en
paliers~:

($\mathord{\boldsymbol{\circ}}$) Prérequis ($\mathord{\bullet}$)
Fondamental ($\mathord{\bullet}\mathord{\bullet}$) Standard
($\mathord{\bullet}\mathord{\bullet}\mathord{\bullet}$) Avancé
($\mathord{\bullet}\mathord{\bullet}\mathord{\bullet}\mathord{\bullet}$)
Expert

Sauf mention particulière, la connaissance des démonstrations du
document n'est pas exigible\footnote{l'étude des démonstrations du cours
  peut toutefois contribuer à votre apprentissage, au même titre que la
  résolution d'exercices.}

\runinhead{Préambule}\phantomsection\label{IV-pruxe9ambule}

\begin{itemize}
\tightlist
\item
  $\mathord{\bullet}$ connaître la définition des espaces
  \(\mathcal{L}^p\)
\item
  $\mathord{\bullet}\mathord{\bullet}$ connaître la propriété
  d'inclusion dans les espaces \(\mathcal{L}^p\)
\item
  $\mathord{\bullet}$ connaître l'inégalité de Markov
\item
  $\mathord{\bullet}$ connaître l'inégalité de Bienaymé-Chebyshev
\end{itemize}

\runinhead{Suites de v.a.
indépendantes}\phantomsection\label{IV-suites-de-v.a.-induxe9pendantes}

\begin{itemize}
\tightlist
\item
  $\mathord{\bullet}\mathord{\bullet}\mathord{\bullet}\mathord{\bullet}$ connaître
  le théorème d'existence et d'unicité de la mesure produit sur les
  espaces produit de dimension infinie
\item
  $\mathord{\bullet}\mathord{\bullet}$ connaître la définition d'une
  suite de variables aléatoires indépendantes
\item
  $\mathord{\bullet}\mathord{\bullet}$ connaître et savoir utiliser le
  lemme de Borel-Cantelli
\end{itemize}

\runinhead{Convergence et loi des grands
nombres}\phantomsection\label{IV-convergence-et-loi-des-grands-nombres}

\begin{itemize}
\tightlist
\item
  $\mathord{\bullet}\mathord{\bullet}$ connaître les différents modes de
  convergence des suites de v.a.
\item
  $\mathord{\bullet}\mathord{\bullet}\mathord{\bullet}$ savoir
  caractériser le(s) mode(s) de convergence d'une suite de v.a.
\item
  $\mathord{\bullet}\mathord{\bullet}$ connaître les liens d'implication
  entre les différents modes de convergence
\item
  $\mathord{\bullet}\mathord{\bullet}$ connaître la propriété de
  continuité
\item
  $\mathord{\bullet}$ connaître la loi faible des grands nombre
\item
  $\mathord{\bullet}$ connaître la loi forte des grands nombre (sous les
  hypothèses minimales)
\end{itemize}

\runinhead{Convergence en loi et théorème central
limite}\phantomsection\label{IV-convergence-en-loi-et-thuxe9oruxe8me-central-limite}

\begin{itemize}
\tightlist
\item
  $\mathord{\bullet}$ connaître la définition de la convergence en loi
\item
  $\mathord{\bullet}$ savoir que la convergence en probabilité implique
  la convergence en loi
\item
  $\mathord{\bullet}\mathord{\bullet}$ savoir que la convergence simple
  des fonctions de répartition équivaut à la convergence en loi
\item
  $\mathord{\bullet}$ connaître et savoir utiliser le théorème central
  limite
\end{itemize}

\chapter{Préambule}\label{IV-pruxe9ambule-1}

Dans ce chapitre, on va s'intéresser à l'étude des suites de variables
aléatoires. On introduira notamment différents modes de convergence. En
particulier, on établira la \emph{loi des grands nombres} qui montre
rigoureusement que, quand le nombre de répétitions de l'expérience tend
vers l'infini, la fréquence de réalisation d'un événement converge vers
la probabilité de réalisation de cet événement. Ce résultat justifie
également la démarche employée en statistique où l'on cherche par
exemple à estimer une espérance par une moyenne empirique~: cette
moyenne empirique va ainsi converger vers la quantité ciblée. Enfin, on
présentera le théorème central limite qui nous indiquera la vitesse à
laquelle cette convergence a lieu ainsi qu'un moyen de contrôler
l'erreur commise en remplaçant une espérance par sa contrepartie
empirique.

Notons d'abord qu'on peut étendre les définitions des espaces
\(\mathcal{L}^1\) et \(\mathcal{L}^2\) pour un \(p \in \mathbb{N}^\ast\)
quelconque.

\addcontentsline{toc}{subsection}{Espace \(\mathcal{L}^p\)}
\begin{definition}[{Espace
\(\mathcal{L}^p\)}]\phantomsection\label{IV-espace-mathcallp}

Soit \(X\) une variable aléatoire. On note \(X \in \mathcal{L}^p\), ou
\(\mathcal{L}^p(\Omega,\mathcal{A},\mathbb{P})\), si et seulement si
\(\mathbb{E}(|X|^p) = \int_\mathbb{R}|x|^p \mathbb{P}_X(dx) = \int_{\Omega} |X|^p(\omega)\mathbb{P}(d\omega) < +\infty\).

Si \(X \in \mathcal{L}^p\), on dit qu'elle admet un moment d'ordre
\(p\). Du fait que \(\mathbb{P}\) est une mesure finie, on a la
stabilité par inclusion suivante~:
\end{definition}

\addcontentsline{toc}{subsection}{Inclusion}
\begin{proposition}[{Inclusion}]\phantomsection\label{IV-inclusion}

Soit \(p \in \mathbb{N}^\ast\), on a l'inclusion~:
\[ \mathcal{L}^{p+1}(\Omega,\mathcal{A},\mathbb{P}) \subset \mathcal{L}^p(\Omega,\mathcal{A},\mathbb{P})\]
\end{proposition}

\begin{preuve}
Supposons \(X \in \mathcal{L}^{p+1}(\Omega,\mathcal{A},\mathbb{P})\). On
a
\[|X|^p \leq \max(1,|X|^{p+1}) = 1_{|X| < 1} + 1_{|X| \geq 1} |X|^{p+1}.\]
Le terme de droite est intégrable, en effet~:
\[\mathbb{E}\left(1_{|X| < 1} + 1_{|X| \geq 1} |X|^{p+1}\right) \leq \int_\Omega \mathbb{P}(d\omega) + \int_\Omega |X(\omega)|^{p+1} \mathbb{P}(d\omega) = 1 + \mathbb{E}(|X|^{p+1}).\]
donc \(|X|^p\) est intégrable.\hfill$\blacksquare$
\end{preuve}

On donne ici dans le cas général les inégalités de Markov et de
Bienaymé-Chebyshev, déjà vues en CPGE.

\addcontentsline{toc}{subsection}{Inégalité de Markov}
\begin{theoreme}[{Inégalité de
Markov}]\phantomsection\label{IV-inegmarkov}

Soient \(p \in \mathbb{N}^\ast\) et \(a \in \mathbb{R}_+^\ast\). Soit
\(X \in \mathcal{L}^p(\Omega,\mathcal{A},\mathbb{P})\), on a
\[\mathbb{P}(|X|\geq a) \leq \frac{\mathbb{E}(|X|^p)}{a^p}\]
\end{theoreme}

\begin{preuve}
On a \[|X|^p \geq a^p 1_{[a, +\infty[}(|X|)\] et en prenant l'espérance,
on obtient ainsi
\[\mathbb{E}(|X|^p) \geq a^p \mathbb{E}(1_{[a, +\infty[}(|X|)) = a^p \mathbb{P}(|X|\geq a).\]

\hfill$\blacksquare$
\end{preuve}

\addcontentsline{toc}{subsection}{Inégalité de Bienaymé-Chebyshev}
\begin{corollaire}[{Inégalité de
Bienaymé-Chebyshev}]\phantomsection\label{IV-inegbc}

Soit \(X \in \mathcal{L}^2\), on a
\[\mathbb{P}(|X-\mathbb{E}(X)| \geq a) \leq \frac{\mathbb{V}(X)}{a^2}\]
\end{corollaire}

\runinhead{Exercice -- Démonstration (\(\mathord{\bullet}\))}\phantomsection\addcontentsline{toc}{subsubsection}{Démonstration}\label{IV-bienaymuxe9cheby}
\solref{IV-answer-bienaymuxe9cheby}

\addcontentsline{toc}{subsection}{Interprétation}
\begin{remarque}[{Interprétation}]\phantomsection\label{IV-interpruxe9tation}

L'inégalité de Bienaymé-Chebyshev est très utile en pratique. Elle
permet de mesurer la probabilité des grands écarts entre \(X\) et sa
moyenne. Par exemple, avec \(a = 10 \sigma\), il en résulte qu'il est
improbable qu'une variable aléatoire \(X\) dévie de son espérance
\(\mathbb{E}(X)\) de plus de 10 fois son écart-type (la probabilité est
inférieure à 0.01). Cette inégalité, tout à fait générale, n'est
cependant pas très précise, et surestime très souvent en pratique le
membre de gauche. On préférera, quand c'est possible, calculer
directement ces probabilités à partir de la loi de \(X\).
\end{remarque}

\chapter{Suites de variables
indépendantes}\label{IV-suites-de-variables-induxe9pendantes}

Au début du chapitre II, nous avons vu comment construire des couples
(et même des \(n\)-uplets en itérant) de variables aléatoires réelles
indépendantes en considérant l'espace produit, munis des tribus produit
et des (mesures de) probabilités produit. Il est malheureusement
beaucoup plus délicat, mais indispensable pour les applications (en
particulier la loi des grands nombres), de construire une \textbf{suite
infinie de variables indépendantes} de lois données.

Plus précisément, pour chaque entier \(n\) on se donne une v.a.r.
\(X_n\) définie sur un espace de probabilité
\((\Omega_n,\mathcal{A}_n,\mathbb{P}_n)\), à valeurs dans
\((\mathbb{R},\mathcal{B}(\mathbb{R}))\) et de loi \(\mathbb{P}_{X_n}\)
(pour construire chaque \(X_n\), on peut procéder comme ci-dessus).
Ensuite, on pose \begin{align*}
\Omega &= \prod_{n=1}^{\infty}\Omega_n \text{        (produit cartésien dénombrable),}\\
\mathcal{A}&= \bigotimes_{n=1}^{\infty}\mathcal{A}_n,
\end{align*} où \(\bigotimes_{n=1}^{\infty}\mathcal{A}_n\) désigne la
plus petite tribu de \(\Omega\) à laquelle appartienne tous les
ensembles de la forme
\[A_1 \times A_2 \times \ldots\times A_k \times \Omega_{k+1} \times \Omega_{k+2} \times \ldots, \,\,\, A_i \in \mathcal{A}_i, \,\,\, k \in \mathbb{N}^\ast\]

On a alors le théorème suivant, que l'on admettra, qui constitue un
résultat non trivial de la théorie de la mesure et qui généralise le
\href{Probabilité\%20III.pdf\%20\#fubiniproba}{théorème de Fubini}.

\addcontentsline{toc}{subsection}{Mesure produit sur un espace de
dimension infinie}
\begin{theoreme}[{Mesure produit sur un espace de dimension
infinie}]\phantomsection\label{IV-probaespaceden}

Avec les notations ci-dessus, il existe une unique probabilité
\(\mathbb{P}\) sur \((\Omega,\mathcal{A})\), telle que
\[\mathbb{P}(A_1 \times \ldots\times A_k \times \Omega_{k+1} \times \ldots) =\prod_{i=1}^k \mathbb{P}_i(A_i)\]
pour tous \(k \in \mathbb{N}^{\ast}\) et \(A_i \in \mathcal{A}_i\),
\(i \in \{1,\ldots,k\}\).

Nous sommes maintenant en mesure de définir une suite de variables
aléatoires toutes définies sur le même espace de probabilité
\((\Omega, \mathcal{A}, \mathbb{P})\).
\end{theoreme}

\addcontentsline{toc}{subsection}{Suites de variables aléatoires
indépendantes}
\begin{definition}[{Suites de variables aléatoires
indépendantes}]\phantomsection\label{IV-suites-de-variables-aluxe9atoires-induxe9pendantes}

La suite \((X_n)_{n\in \mathbb{N}^\ast}\) de variables aléatoires est
dite \emph{indépendante} si pour tout \(n \in \mathbb{N}^\ast\), la
famille finie \(X_1,\ldots,X_n\) est indépendante.

Il est facile de vérifier que l'indépendance est préservée par certaines
transformations.
\end{definition}

\addcontentsline{toc}{subsection}{Conséquences}
\begin{proposition}[{Conséquences}]\phantomsection\label{IV-consuxe9quences}

L'indépendance de la suite \((X_n)_{n\in \mathbb{N}^\ast}\) entraîne
celle de

\begin{enumerate}
\def\labelenumi{\arabic{enumi}.}
\tightlist
\item
  toute sous-suite \((X_{i_k})_{k\in \mathbb{N}^\ast}\),
\item
  toute suite de vecteurs issus de \(X_n\),
\item
  toute suite de la forme \((f_n(X_n))_{n\in \mathbb{N}^\ast}\), où les
  fonctions \(f_n\) sont des fonctions mesurables.
\end{enumerate}
\end{proposition}

\addcontentsline{toc}{subsection}{Développement dyadique}
\begin{exemple}[{Développement
dyadique}]\phantomsection\label{IV-duxe9veloppement-dyadique}

Nous considérons l'ensemble \(\Omega = [0, 1[\) muni de la tribu
borélienne restreinte à cet ensemble, et de la mesure de Lebesgue. A
chaque réel \(\omega\), nous associons son développement dyadique
(unique si l'on impose que les \(\omega_i\) ne soient pas tous égaux à 1
à partir d'un certain rang)~:
\[ \omega = \sum_{i\in \mathbb{N}^\ast} \frac{\omega_i}{2^i},\,\,\, \omega_i \in \{0,1\}.\]
L'application \(X_i : \Omega \to \{0,1\}\), qui à \(\omega\) associe
\(X_i(\omega) = \omega_i\) est une variable aléatoire sur \(\Omega\). En
effet, pour \(x_i \in \{0,1\}\), \(1\leq i \leq n\),
\[ \{X_i = x_i\} = \bigcup_{x_1,\ldots, x_{i-1} \in \{0,1\}} \left[ \sum_{j=1}^i \frac{x_j}{2^j}, \sum_{j=1}^i \frac{x_j}{2^j} + \frac{1}{2^i}\right[, \]
qui est bien un élément de la tribu borélienne de \(\Omega = [0,1[\), et
\[\mathbb{P}(\{X_i = x_i\}) = \frac{1}{2^i} \sum_{x_1,\ldots, x_{i-1} \in \{0,1\}} 1 = \frac{1}{2}.\]
Montrons l'indépendance des variables aléatoires
\((X_i)_{1\leq i \leq n}\). Nous avons
\[ \bigcap_{1\leq i \leq n} \{X_i = x_i\} = \left[ \sum_{i=1}^n \frac{x_i}{2^i}, \sum_{i=1}^n \frac{x_i}{2^i} + \frac{1}{2^n}\right[ \]
si bien que
\[ \mathbb{P}\left(\bigcap_{1\leq i \leq n} \{X_i = x_i\}\right) = \frac{1}{2^n} = \prod_{1\leq i \leq n} \mathbb{P}(\{X_i = x_i\}).\]
Cela démontre que les variables aléatoires \(X_i\) sont indépendantes et
de loi de Bernoulli de paramètre \(\frac{1}{2}\) . Ainsi, nous venons de
construire sur \(\Omega\) une suite de variables aléatoires telle que
toute sous-suite finie est constituée de variables aléatoires
indépendantes. C'est donc une suite de variables aléatoires
indépendantes de loi de Bernoulli de paramètre \(\frac{1}{2}\) .

On note enfin le résultat très utile suivant portant sur les suites
d'événements indépendants.
\end{exemple}

\addcontentsline{toc}{subsection}{Borel-Cantelli}
\begin{lemme}[{Borel-Cantelli}]\phantomsection\label{IV-BC}

Soit \((A_n)_{n\in\mathbb{N}^\ast}\) une suite d'événements sur l'espace
probabilisé \((\Omega,\mathcal{A}, \mathbb{P})\).

\begin{enumerate}
\def\labelenumi{\arabic{enumi}.}
\item
  Si \(\sum_{n=1}^{\infty} \mathbb{P}(A_n) < \infty\), alors
  \(\mathbb{P}(\limsup\limits_{n \to \infty} A_n) = \mathbb{P}\left(\bigcap_{n\geq 1} \bigcup_{k \geq n} A_k \right) =0\).
\item
  Si \(\sum_{n=1}^{\infty} \mathbb{P}(A_n) = \infty\) et si les
  événements \(A_n\) sont mutuellement indépendants, alors on a
  \(\mathbb{P}(\limsup\limits_{n \to \infty} A_n) = 1\).
\end{enumerate}
\end{lemme}

\begin{preuve}
Exercice.\hfill$\blacksquare$
\end{preuve}

\runinhead{Exercice -- Pile ou face
infini (\(\mathord{\bullet}\))}\phantomsection\addcontentsline{toc}{subsubsection}{Pile
ou face infini}\label{IV-pf}
On considère un jeu de pile ou face infini où on a une probabilité
\(p, 0< p < 1\), d'obtenir face. Montrer que l'événement
\(A = \{\omega : \text{ il n'y a qu'un nombre fini de faces}\}\) est de
probabilité nulle. (Considérer les événements
\(A_n = \{\text{on a face au }n\text{-ième tirage}\}\), puis montrer que
\(\sum_{n=1}^\infty \mathbb{P}(A_n) = \infty\)).\solref{IV-answer-pf}

\chapter{Convergences et loi des grands
nombres}\label{IV-convergences-et-loi-des-grands-nombres}

\section{Convergences des variables
aléatoires}\label{IV-convergences-des-variables-aluxe9atoires}

Dans ce paragraphe, on va décrire les notions de convergence de
variables (on y inclut les vecteurs) aléatoires. On verra que plusieurs
notions sont possibles, non équivalentes, ce qui enrichit mais complique
aussi la description des comportements asymptotiques.

En calcul intégral, on a beaucoup étudié le cas de suites de fonctions
convergeant simplement. Une variable aléatoire étant une fonction, on a
donc la même notion qui revient à écrire qu'une suite
\((X_n)_{n\in \mathbb{N}^\ast}\) de variables aléatoires converge
simplement vers \(X\) si \(\lim_{n \to \infty} X_n(\omega) = X(\omega)\)
pour tout \(\omega \in \Omega\). Cette définition naturelle est
malheureusement à peu près inutile en probabilités, comme l'illustre
l'exemple suivant.

\addcontentsline{toc}{subsection}{Pile ou face infini (bis)}
\begin{exemple}[{Pile ou face infini
(bis)}]\phantomsection\label{IV-pfinf}

Soit \((X_n)_{n\in \mathbb{N}^\ast}\) une suite de v.a. réelles qui sont
indépendantes et de même loi (on notera \emph{i.i.d} pour
\emph{indépendantes et identiquement distribuées}), avec
\(\mathbb{P}(X_n = 1) = p\) et \(\mathbb{P}(X_n=0) = 1-p\). Ce sont donc
des v.a. de loi de Bernoulli de paramètre \(p\) qui modélisent par
exemple les résultats d'un jeu de pile ou face. Lorsque \(n\) est grand,
on s'attend à ce que la proportion de faces (\(X_n=1\)) soit à peu près
égale à \(p\) (c'est l'essence de la conception objectiviste des
probabilités). Mathématiquement, on voudrait que
\[\lim_{n \to \infty} \frac{X_1(\omega) + \ldots +X_n(\omega)}{n} = p \text{   pour tout }\omega\in\Omega.\]
C'est en fait complètement faux. En effet, si on considère par exemple
la suite \(\omega_0 = {p,p,p,\ldots}\) qui ne contient que des piles, on
obtient
\[\lim_{n \to \infty} \frac{X_1(\omega_0) + \ldots +X_n(\omega_0)}{n} = 0\]
et plus généralement, on a
\(\lim_{n \to \infty} \frac{1}{n}\sum_{i=1}^n X_i(\omega) = 0\) pour
tout \(\omega\) dans l'ensemble
\(A = \{\omega : \text{ il n'y a qu'un nombre fini de faces}\}\). On
peut même trouver des \(\omega\) pour lesquels la fréquence converge
vers n'importe quel nombre fixé dans \([0,1]\)\,! Bien entendu,
l'événement \(A\) est invraisemblable et on peut montrer qu'il vérifie
\(\mathbb{P}(A) = 0\)\footnote{cf.~\protect\hyperref[IV-pf]{exercice
  plus haut (p.~\pageref*{IV-pf})}}. On verra cependant que la
\textbf{loi des grands nombres} assure~:
\[\mathbb{P}\left(\left\{\omega ; \lim_{n \to \infty} \frac{1}{n}\sum_{i=1}^n X_i(\omega) = p\right\}\right)=1.\]
Ce type de convergence, pour lequel on n'a pas convergence pour tout
\(\omega\), mais seulement pour presque tout \(\omega\) (autrement dit
\(\mathbb{P}\)-presque partout) est typiquement ce qui arrive pour des
variables aléatoires.

On considère une suite \((X_n)_{n\in \mathbb{N}^\ast}\) de variables
aléatoires, définies sur un même espace de probabilité
\((\Omega, \mathcal{A}, \mathbb{P})\), et à valeurs dans
\(\mathbb{R}^d\). On considère également sur le même espace un vecteur
``limite'' \(X\). On notera \(|\cdot|\) la valeur absolue dans
\(\mathbb{R}\) ou la norme euclidienne dans \(\mathbb{R}^d\).
\end{exemple}

\addcontentsline{toc}{subsection}{Modes de convergences}
\begin{definition}[{Modes de
convergences}]\phantomsection\label{IV-modes-de-convergences}

\begin{enumerate}
\def\labelenumi{\arabic{enumi}.}
\tightlist
\item
  La suite \((X_n)_{n\in \mathbb{N}^\ast}\) \emph{converge presque
  sûrement} vers \(X\), ce qui s'écrit \(X_n \to X\) p.s., s'il existe
  un ensemble \(N \in \mathcal{A}\) de probabilité nulle tel que
  \[ X_n (\omega) \xrightarrow[n \to \infty]{}  X(\omega),\,\,\, \forall \omega \notin N.\]
\item
  La suite \((X_n)_{n\in \mathbb{N}^\ast}\) \emph{converge en
  probabilité} vers \(X\), ce qui s'écrit
  \(X_n \xrightarrow{\mathbb{P}} X\), si \(\forall \varepsilon >0\) on a
  \[ \mathbb{P}(|X_n-X| \geq \varepsilon) \xrightarrow[n \to \infty]{}  0. \]
\item
  La suite \((X_n)_{n\in \mathbb{N}^\ast}\) \emph{converge en moyenne}
  (ou dans \(\mathcal{L}^1\)) vers \(X\), ce qui s'écrit
  \(X_n \xrightarrow{\mathcal{L}^1} X\), si \(X_n\) et \(X\) sont dans
  \(\mathcal{L}^1\) et si
  \[ \mathbb{E}(|X_n - X|) \xrightarrow[n \to \infty]{}  0 .\]
\end{enumerate}
\end{definition}

\addcontentsline{toc}{subsection}{Convergence \(\mathcal{L}^p\)}
\begin{remarque}[{Convergence
\(\mathcal{L}^p\)}]\phantomsection\label{IV-convergence-mathcallp}

La définition de la convergence dans \(\mathcal{L}^1\) se généralise aux
ordres supérieurs, pour \(p \in \mathbb{N}^\ast\), on parle de
convergence dans \(\mathcal{L}^p\) (en moyenne quadratique si \(p=2\))
ce qui s'écrit \(X_n \xrightarrow{\mathcal{L}^p} X\), si \(X_n\) et
\(X\) sont dans \(\mathcal{L}^p\) et si
\[ \mathbb{E}(|X_n - X|^p) \xrightarrow[n \to \infty]{}  0 .\]

Ces modes de convergences ne sont pas équivalents comme le montrent les
exemples suivants.
\end{remarque}

\runinhead{Exercice -- En proba et en
moyenne (\(\mathord{\bullet}\))}\phantomsection\addcontentsline{toc}{subsubsection}{En
proba et en moyenne}\label{IV-lien}
Soit \((X_n)_{n\in \mathbb{N}^\ast}\) une suite de variables aléatoires
de Bernoulli à valeurs dans \(\{0,1\}\) telles que
\[ \mathbb{P}(X_n=1) = \frac{1}{n} ;\,\,\,\, \mathbb{P}(X_n = 0) = 1- \frac{1}{n}.\]
Montrer que la suite \((X_n)_{n\in \mathbb{N}^\ast}\) converge en
probabilité et en moyenne vers 0.\solref{IV-answer-lien}

\runinhead{Exercice -- En proba mais pas en
moyenne (\(\mathord{\bullet}\))}\phantomsection\addcontentsline{toc}{subsubsection}{En
proba mais pas en moyenne}\label{IV-lien2}
Soit maintenant une suite \((Y_n)_{n\in \mathbb{N}^\ast}\) de variables
aléatoires de Bernoulli à valeurs dans \(\{0,n^2\}\) telles que
\[ \mathbb{P}(Y_n=n^2) = \frac{1}{n} ;\,\,\,\, \mathbb{P}(Y_n = 0) = 1- \frac{1}{n}.\]
Montrer que la suite \((Y_n)_{n\in \mathbb{N}^\ast}\) converge en
probabilité vers 0 mais pas en moyenne.\solref{IV-answer-lien2}

\runinhead{Exercice -- Presque
sûre (\(\mathord{\bullet}\))}\phantomsection\addcontentsline{toc}{subsubsection}{Presque
sûre}\label{IV-lien3}
Soit \(U\) une variable aléatoire uniforme sur \([0 , 1]\). Posons
\(Z_n = 1_{\{ U \leq \frac{1}{n}\} }\). Alors
\[ \mathbb{P}(Z_n = 1)= \frac{1}{n} ;\,\,\,\, \mathbb{P}(Z_n = 0) = 1 - \frac{1}{n}\]
Montrer que la suite \((Z_n)_{n\in\mathbb{N}^\ast}\) converge presque
sûrement vers 0.\solref{IV-answer-lien3}

Il existe cependant des liens entre ces différents modes de convergences
que l'on explore ci-dessous.

\addcontentsline{toc}{subsection}{p.s. / \(\mathcal{L}^1\)
\(\Rightarrow\) en proba}
\begin{proposition}[{p.s. / \(\mathcal{L}^1\) \(\Rightarrow\) en
proba}]\phantomsection\label{IV-propconv1}

La convergence presque sûre et la convergence en moyenne entraînent la
convergence en probabilité.
\end{proposition}

\begin{preuve}
Soient \(n\in\mathbb{N}^\ast\), \(\varepsilon > 0\) et
\(A_{n,\varepsilon} = \{|X_n - X| > \varepsilon \}\).

\begin{itemize}
\item
  Supposons que \(X_n \to X\) p.s. et soit \(N\) l'ensemble de
  probabilité nulle en dehors duquel on a \(X_n(\omega) \to X(\omega)\).
  Si \(\omega \notin N\), on a \(\omega \notin A_{n,\varepsilon}\) pour
  tout \(n \geq n_0\), où \(n_0\) dépend de \(\omega\) et de
  \(\varepsilon\), ce qui implique que les variables aléatoires
  \(Y_{n,\varepsilon} = 1_{N^c\cap A_{n,\varepsilon}}\) tendent
  simplement vers 0 lorsque \(n \to \infty\). Comme on a aussi
  \(0 \leq Y_{n,\varepsilon} \leq 1\) le théorème de convergence dominée
  implique que
  \(\mathbb{E}(Y_{n,\varepsilon}) \xrightarrow[n \to \infty]{} 0.\) Par
  suite,
  \[\mathbb{P}(A_{n,\varepsilon}) \leq \mathbb{P}(N^c \cap A_{n,\varepsilon}) + \mathbb{P}(N) = \mathbb{P}(N^c \cap A_{n,\varepsilon}) = \mathbb{E}(Y_{n,\varepsilon}) \xrightarrow[n \to \infty]{} 0.\]
\item
  Supposons que \(X_n \xrightarrow{\mathcal{L}^1} X\). Pour
  \(\varepsilon >0\), on a
  \(1_{A_{n,\varepsilon}} \leq \frac{1}{\varepsilon}|X_n - X|\), donc
  \[ \mathbb{P}(A_{n,\varepsilon}) \leq \frac{1}{\varepsilon}\mathbb{E}(|X_n - X|) \xrightarrow[n \to \infty]{} 0.\]
\end{itemize}

\hfill$\blacksquare$
\end{preuve}

La convergence en probabilité n'entraîne pas la convergence en moyenne,
comme nous l'avons vu dans \protect\hyperref[IV-lien2]{l'exemple
précédent (p.~\pageref*{IV-lien2})}, ne serait-ce que parce qu'elle
n'implique pas l'appartenance de \(X_n\) et \(X\) à \(\mathcal{L}^1\).
Si les \(X_n\) ne sont ``pas trop grands'', il y a cependant équivalence
entre les deux modes de convergence. En voici un exemple.

\addcontentsline{toc}{subsection}{Cas borné}
\begin{proposition}[{Cas borné}]\phantomsection\label{IV-propconv2}

S'il existe une constante \(a\) telle que \(|X_n| \leq a\) presque
sûrement, il y a équivalence entre \(X_n \xrightarrow{\mathbb{P}} X\) et
\(X_n \xrightarrow{\mathcal{L}^1} X\).
\end{proposition}

\begin{preuve}
Etant donnée la \protect\hyperref[IV-propconv1]{proposition
précédente (p.~\pageref*{IV-propconv1})}, dont on reprend les notations,
il suffit de montrer que la convergence en probabilité implique la
convergence en moyenne lorsque \(|X_n| \leq a\).

Comme \(|X_n| \leq a\) p.s., on a
\(\{|X| > a +\varepsilon\} \subset A_{n,\varepsilon}\), et donc
\(\mathbb{P}(|X| > a +\varepsilon ) \leq \mathbb{P}(A_{n,\varepsilon})\).
En faisant \(n \to \infty\), on en déduit que
\(\mathbb{P}(|X| > a + \varepsilon) = 0\). Ceci est vrai pour tout
\(\varepsilon\) et donc \[\mathbb{P}(|X|> a) = 0.\] Comme
\(|X_n| \leq a\), on a aussi
\[ |X_n -X| \leq \varepsilon + (X_n + X) 1_{A_{n,\varepsilon}} \leq \varepsilon + 2a 1_{A_{n,\varepsilon}}\]
sur l'ensemble \(\{|X| \leq a \}\) qui est de probabilité 1. On a ainsi
\[ \mathbb{E}(|X_n - X|) \leq \varepsilon +2a \mathbb{P}(A_{n,\varepsilon})\]
On en déduit que
\(\lim \sup_{n \to \infty} \mathbb{E}(|X_n - X|) \leq \varepsilon\), et
comme \(\varepsilon\) est arbitrairement proche de 0, on a le résultat
souhaité.\hfill$\blacksquare$
\end{preuve}

Les rapports entre convergence presque-sûre et convergence en
probabilité sont plus subtils. La première de ces deux convergences est
plus forte que la seconde d'après la
\protect\hyperref[IV-propconv1]{proposition plus
haut (p.~\pageref*{IV-propconv1})}, mais ``à peine plus'', comme le
montre le résultat suivant.

\addcontentsline{toc}{subsection}{Existence d'une sous-suite
convergente}
\begin{proposition}[{Existence d'une sous-suite
convergente}]\phantomsection\label{IV-propconv3}

Si \(X_n \xrightarrow{\mathbb{P}} X\), il existe une sous-suite
\((n_k)\) telle que \(X_{n_k} \to X\) p.s. quand \(k \to \infty\).
\end{proposition}

\begin{preuve}
On remarque d'abord que l'on peut réécrire la définition de la
convergence en probabilité de la manière suivante~:
\(\forall \varepsilon >0\) et \(\delta > 0\), il existe
\(N = N(\delta,\varepsilon)\) tel que
\[n\geq N \Rightarrow \mathbb{P}(|X_n-X| > \varepsilon ) < \delta.\]

Comme la suite \((X_n)_{n\in \mathbb{N}^\ast}\) converge en probabilité
vers \(X\), on peut définir une sous-suite de la manière suivante~:
posons \(n_1 = 1\), et soit
\[n_j = \inf \left\{ n > n_{j-1} ; \mathbb{P}\left(|X_m - X| > \frac{1}{2^j} \right) < \frac{1}{3^j}, \text{ pour } m \geq n \right\}.\]
On a alors
\[\sum_j \mathbb{P}\left(|X_{n_j} - X| > \frac{1}{2^j} \right) < \sum_j \frac{1}{3^j} < \infty\]
et en appliquant le \protect\hyperref[IV-BC]{lemme de
Borel-Cantelli (p.~\pageref*{IV-BC})} aux ensembles
\[A_j =\left\{|X_{n_{j}} - X| > \frac{1}{2^j} \right\},\] on obtient que
la suite \((X_{n_j})_{j\in \mathbb{N}^\ast}\) converge
presque-sûrement.\hfill$\blacksquare$
\end{preuve}

\runinhead{Exercice -- Intervalles de longueur aléatoire
(1) (\(\mathord{\bullet}\))}\phantomsection\addcontentsline{toc}{subsubsection}{Intervalles
de longueur aléatoire (1)}\label{IV-ia}
Soient \(\Omega = \mathbb{R}\) muni de sa tribu borélienne et
\(\mathbb{P}\) la probabilité uniforme sur {[}0, 1{]}. Soit
\(X_n = 1_{A_n}\) , où \(A_n\) est un intervalle de {[}0, 1{]} de
longueur \(\frac{1}{n}\).

Montrer que la suite des \(X_n\) converge en moyenne et en
proba.\solref{IV-answer-ia}

\runinhead{Exercice -- Intervalles de longueur aléatoire
(2) (\(\mathord{\bullet}\mathord{\bullet}\mathord{\bullet}\))}\phantomsection\addcontentsline{toc}{subsubsection}{Intervalles
de longueur aléatoire (2)}\label{IV-ia2}
Montrer que la suite des \(X_n\) ne converge pas presque sûrement.
Extraire une sous-suite qui converge presque
sûrement.\solref{IV-answer-ia2}

\addcontentsline{toc}{subsection}{Continuité}
\begin{proposition}[{Continuité}]\phantomsection\label{IV-propconv4}

Soit \(f\) une fonction continue de \(\mathbb{R}^d\) dans
\(\mathbb{R}\).

\begin{enumerate}
\def\labelenumi{\arabic{enumi}.}
\tightlist
\item
  Si \(X_n \to X\) presque-sûrement, alors \(f(X_n) \to f(X)\)
  presque-sûrement.
\item
  Si \(X_n \xrightarrow{\mathbb{P}} X\), alors
  \(f(X_n) \xrightarrow{\mathbb{P}} f(X)\).
\end{enumerate}
\end{proposition}

\begin{preuve}
\begin{enumerate}
\def\labelenumi{\arabic{enumi}.}
\item
  Soit \(N\) l'ensemble de probabilité nulle en dehors duquel on a
  \(X_n(\omega) \xrightarrow[n \to \infty]{} X(\omega)\). Si
  \(\omega \notin N\), il vient
  \[ \lim_{n \to \infty}f(X_n(\omega)) = f(\lim_{n \to \infty}X_n(\omega)) = f(X(\omega))\]
  par continuité de \(f\), d'où le résultat.
\item
  On remarque d'abord que si \(K >0\) et \(\varepsilon >0\),
  \[\{|f(X_n)-f(X)| \geq \varepsilon\} \subset \{|X| > K\}\cup\{|X| \leq K, |f(X_n)-f(X)| \geq \varepsilon\}.\]
  La fonction \(f\) est uniformément continue sur
  \(\{x : |x| \leq K\}\), donc il existe \(\eta > 0\) tel que
  \(|x-y| <\eta\) et \(|x| \leq 2K\) et \(|y|\leq 2K\) impliquent
  \(|f(x) - f(y)| <\varepsilon\). On a donc
  \[\{|f(X_n)-f(X)| \geq \varepsilon\} \subset \{|X| > K\}\cup\{|X_n-X| \geq \eta\}\]
  d'où
  \[\mathbb{P}(|f(X_n)-f(X)| \geq \varepsilon) \leq \mathbb{P}(|X| > K) + \mathbb{P}(|X_n-X| \geq \eta).\]
  Par hypothèse, il vient
  \[ \lim_{n \to \infty} \mathbb{P}(|f(X_n)-f(X)| \geq \varepsilon) \leq \mathbb{P}(|X| > K) .\]
  Enfin, \(\lim_{K\to +\infty} \mathbb{P}(|X| > K) = 0\) (par
  convergence dominée) et donc la limite ci-dessus est nulle. On a ainsi
  le résultat.
\end{enumerate}

\hfill$\blacksquare$
\end{preuve}

\section{La loi des grands nombres}\label{IV-la-loi-des-grands-nombres}

On présente maintenant l'un des résultats essentiels de la théorie des
probabilités. Ce résultat montre rigoureusement que, quand le nombre de
répétitions de l'expérience tend vers l'infini, la fréquence de
réalisation d'un événement converge p.s. vers la probabilité de
réalisation de cet événement. Ce résultat, appelé \textbf{Loi des grands
nombres}, a d'autres portées fondamentales. Il est en particulier à
l'origine de méthodes de calcul numérique appelées Méthodes de
Monte-Carlo, qui sont extrêmement puissantes et robustes. Elles sont par
exemple très utilisées en Physique, en Mathématiques Financières ou
encore dans les problèmes de quantification d'incertitudes. Une
introduction à ces méthodes sera donnée au chapitre V.

Dans ce paragraphe, on considère une suite
\((X_n)_{n\in\mathbb{N}^\ast}\) de variables aléatoires
\textbf{indépendantes et de même loi} (ou indépendantes et identiquement
distribuées, i.i.d. en abrégé). On considère la ``moyenne'' des \(n\)
premières variables aléatoires~:
\[ M_n = \frac{X_1 + \ldots + X_n}{n},\] et notre objectif est de
montrer que \(M_n\) converge vers l'espérance des \(X_n\) lorsque
celle-ci existe (comme les \(X_n\) ont même loi, cette espérance est la
même pour tout \(n\)).

Nous allons démontrer dans un premier temps la loi des grands nombres
pour des variables aléatoires de carré intégrable.

\addcontentsline{toc}{subsection}{Loi des grands nombres cas
\(\mathcal{L}^2\)}
\begin{theoreme}[{Loi des grands nombres cas
\(\mathcal{L}^2\)}]\phantomsection\label{IV-loi-des-grands-nombres-cas-mathcall2}

Soit \((X_n)_{n\in\mathbb{N}^\ast}\) une suite de variables aléatoires
indépendantes, de même loi et de \textbf{carré intégrable}, et
\(m = \mathbb{E}(X_n)\) leur espérance. Alors la suite
\((M_n)_{n\in\mathbb{N}^\ast}\) définie par
\[M_n = \frac{X_1 + \ldots + X_n}{n}\] converge vers \(m\),
\textbf{presque sûrement et en moyenne}, quand \(n\) tend vers l'infini.
Elle converge donc aussi en probabilité. On a même convergence en
\emph{moyenne quadratique}, à savoir que~:
\[ \mathbb{E}((M_n - m)^2) \xrightarrow[n \to \infty]{} 0.\]
\end{theoreme}

\runinhead{Exercice -- Loi faible des grands
nombres (\(\mathord{\bullet}\))}\phantomsection\addcontentsline{toc}{subsubsection}{Loi
faible des grands nombres}\label{IV-wlln}
Le résultat sur la convergence en probabilité est appelé \emph{loi
faible des grands nombres}. Démontrer ce
résultat.\solref{IV-answer-wlln}

Le résultat est peu informatif et permet d'obtenir certains contrôles
d'erreurs. Le résultat prouvant la convergence presque-sûre est appelé
\emph{loi forte des grands nombres}. Sa preuve est plus délicate et
utilise le \protect\hyperref[IV-BC]{lemme de
Borel-Cantelli (p.~\pageref*{IV-BC})}.

\subsection{Démonstration}\label{IV-duxe9monstration-7}

Notons \(\sigma^2\) la variance des variables \(X_n\), bien définie
puisqu'on les a supposées de carré intégrable. En vertu de la linéarité
de l'espérance, on a
\[ \mathbb{E}(M_n) = m,\,\,\,\, \mathbb{E}((M_n-m)^2) = \mathbb{V}(M_n) = \frac{\sigma^2}{n},\]
d'où la convergence en moyenne quadratique.

Comme \(\mathbb{E}(Y)^2 \leq \mathbb{E}(Y^2)\), on en déduit que
\(\mathbb{E}(|M_n -m|)\to 0\), donc on a aussi la convergence en
moyenne.

La preuve de la convergence presque-sûre est plus délicate.

Quitte à remplacer \(X_n\) par \(X_n - m\) (et donc \(M_n\) par
\(M_n - m\)), nous pouvons supposer que \(m = 0\).

Montrons tout d'abord que la sous-suite
\((M_{n^2})_{n \in \mathbb{N}^\ast}\) converge presque-sûrement vers 0.

La convergence dans \(\mathcal{L}^1\) impliquant la convergence en
probabilité, on sait qu'on peut extraire de \((M_n)_n\) une sous-suite
convergeant p.s. vers 0. Cependant, cela ne suffit pas puisqu'on veut
que la suite \((M_n)_n\) elle-même converge p.s. Pour le montrer, on
construit d'abord une sous suite-particulière qui converge p.s. puis on
traite les termes qui se trouvent entre deux éléments successifs de la
sous-suite.

D'après l'\protect\hyperref[IV-inegbc]{inégalité de
Bienaymé-Chebyshev (p.~\pageref*{IV-inegbc})} et ce qui précède, on a
pour \(q\in \mathbb{N}^\ast\)
\[\mathbb{P}\left(|M_{n^2}|\geq \frac{1}{q}\right) \leq \frac{\sigma^2 q^2}{n^2}\]

Donc si \(A_{n,q} = \{|M_{n^2}|\geq \frac{1}{q}\}\), nous obtenons que
\(\sum_{n\geq 1} \mathbb{P}(A_{n,q}) < \infty\). Posons ensuite
\(B_{n,q} = \cup_{m\geq n} A_{m,q}\) et
\(C_q = \cap_{n \geq 1} B_{n,q} = \lim\sup_n A_{n,q}\). En appliquant le
\protect\hyperref[IV-BC]{lemme de Borel-Cantelli (p.~\pageref*{IV-BC})},
on obtient que \(\mathbb{P}(C_q) = 0\). En conséquence, si on pose
\(N = \cup_{q \in \mathbb{N}^\ast} C_q\), on obtient
\(\mathbb{P}(N) \leq \sum_{q\in\mathbb{N}^\ast} \mathbb{P}(C_q)= 0\).

Si \(\omega \notin N\), alors
\(\omega \in \cap_{q \in \mathbb{N}^\ast} (C_q)^c\). Ainsi,
\(\omega \notin C_q\) pour tout \(q \geq 1\), et donc
\(\omega \notin B_{n,q}\) pour \(n\) assez grand (car \(B_{n,q}\) est
décroissant en \(n\)). Cela signifie que pour tout \(\omega \notin N\),
pour tout \(q \geq 1\), il existe un \(n\) assez grand tel que
\(M_{k^2} \leq \frac{1}{q}\) dès que \(k \geq n\). Autrement dit,
\(M_{n^2} \to 0\) si \(\omega \notin N\), avec \(\mathbb{P}(N) = 0\),
d'où \[ M_{n^2} \xrightarrow[n \to \infty]{} 0 \text{ p.s.}\]

Montrons maintenant que la suite \((M_n)_{n\in\mathbb{N}^\ast}\) tend
presque-sûrement vers 0.

Pour tout entier \(n\), notons \(p(n)\), l'entier tel que
\(p(n)^2 \leq n \leq (p(n)+1)^2\). Alors,
\[ M_n - \frac{p(n)^2}{n}M_{p(n)^2} = \frac{1}{n} \sum_{p = p(n)^2+1}^{n} X_p,\]
et puisque les variables aléatoires de la somme sont indépendantes, il
vient \begin{align*}
\mathbb{E}\left(\left(M_n - \frac{p(n)^2}{n}M_{p(n)^2}\right)^2\right) & = \frac{n-p(n)^2}{n^2}\sigma^2 \\
                                                                 & \leq \frac{2p(n)+1}{n^2} \sigma^2 \\
                                                                 & \leq \frac{2\sqrt{n}+1}{n^2} \sigma^2 \leq \frac{3}{n^{3/2}}\sigma^2
\end{align*}

où on a utilisé le fait que \(p(n) \leq \sqrt{n}\).

En appliquant de nouveau l'\protect\hyperref[IV-inegbc]{inégalité de
Bienaymé-Chebyshev (p.~\pageref*{IV-inegbc})}, on obtient
\[ \mathbb{P}\left(\left|M_n - \frac{p(n)^2}{n}M_{p(n)^2}\right|>a\right) \leq \frac{2\sqrt{n}+1}{n^2} \frac{\sigma^2}{a^2}\]
Comme la série \(\sum_n \frac{2\sqrt{n}+1}{n^2}\) converge, le même
raisonnement que pour \(M_{n^2}\) décrit ci-dessus, montre que
\[ M_n - \frac{p(n)^2}{n}M_{p(n)^2} \to 0 \text{ p.s.}\] Par ailleurs,
on a déjà montré que \(M_{p(n)^2} \to 0\) p.s. et
\(\frac{p(n)^2}{n} \to 1\). On en déduit que \(M_n \to 0\) p.s.

Plus généralement, le théorème suivant donne les hypothèses minimales
assurant la validité de la loi des grands nombres, à savoir que les
\(X_n\) sont dans \(\mathcal{L}_1\) (on se référera par exemple à
\href{https://perso.univ-rennes1.fr/ismael.bailleul/AGREG/COURS/LFGN.pdf}{ce
document en ligne} ou à Jacod and Protter
(\protect\hyperlink{refIV-Jacod}{2003}) pour la démonstration).

\addcontentsline{toc}{subsection}{Loi des grands nombres cas
\(\mathcal{L}^1\)}
\begin{theoreme}[{Loi des grands nombres cas
\(\mathcal{L}^1\)}]\phantomsection\label{IV-lfgn}

Soit \((X_n)_{n\in\mathbb{N}^\ast}\) une suite de variables aléatoires
indépendantes, de même loi et \textbf{intégrables}, et
\(m = \mathbb{E}(X_n)\) leur espérance. Alors la suite
\((M_n)_{n\in\mathbb{N}^\ast}\) définie par
\[M_n = \frac{X_1 + \ldots + X_n}{n}\] converge vers \(m\),
\textbf{presque sûrement et en moyenne}, quand \(n\) tend vers l'infini.
\end{theoreme}

\runinhead{Exercice -- Exponentielle (\(\mathord{\bullet}\))}\phantomsection\addcontentsline{toc}{subsubsection}{Exponentielle}\label{IV-expo}
Soit \((X_i)_{i \in \mathbb{N}^\ast}\) une suite de variables aléatoires
i.i.d. d'espérance \(m\) et \(Y_i = e^{X_i}\). Montrer que~:
\[\left(\prod_{i=1}^n Y_i\right)^{1/n}\] converge presque sûrement vers
une constante à déterminer.\solref{IV-answer-expo}

\runinhead{Exercice -- Variance (\(\mathord{\bullet}\))}\phantomsection\addcontentsline{toc}{subsubsection}{Variance}\label{IV-variance}
Soit \((X_i)_{i \in \mathbb{N}^\ast}\) une suite de variables aléatoires
i.i.d. d'espérance \(m\) et de variance \(\sigma^2\). Montrer que~:
\[\lim_{n \to \infty} \frac{1}{n}\sum_{i=1}^n (X_i - m)^2 \to \sigma^2 \text{   p.s.}\]

\solref{IV-answer-variance}

\chapter{Convergence en loi et théorème central
limite}\label{IV-convergence-en-loi-et-thuxe9oruxe8me-central-limite-1}

Nous allons introduire maintenant une nouvelle notion de convergence de
suites de variables aléatoires. La convergence en loi définie dans ce
paragraphe va concerner les lois des variables aléatoires. Elle
signifiera que les lois sont asymptotiquement ``proches'', sans que les
variables aléatoires elles-mêmes le soient nécessairement.

\section{Convergence en loi}\label{IV-convergence-en-loi}

On considère des vecteurs aléatoires \(X_n\) et \(X\), tous à valeurs
dans le même espace \(\mathbb{R}^d\), mais pouvant éventuellement être
définis sur des espaces de probabilité différents.

\addcontentsline{toc}{subsection}{Convergence en loi}
\begin{definition}[{Convergence en
loi}]\phantomsection\label{IV-defconvloi}

On dit que la suite \((X_n)_{n\in \mathbb{N}^\ast}\) \emph{converge en
loi} vers \(X\) et on écrit \(X_n \xrightarrow{\mathcal{L}} X\), si pour
toute fonction réelle \(f\) \textbf{continue bornée} sur
\(\mathbb{R}^d\),
\[\mathbb{E}(f(X_n)) \xrightarrow[n \to \infty]{} \mathbb{E}(f(X)).\]
\end{definition}

\addcontentsline{toc}{subsection}{Cas fini}
\begin{exemple}[{Cas fini}]\phantomsection\label{IV-cas-fini}

Un cas très simple est celui où toutes les variables aléatoires \(X_n\)
prennent un nombre fini de valeurs \(\{ x_i , 1 \leq i \leq N \}\).
Alors, la suite \((X_n)_{n \in \mathbb{N}^\ast}\) converge en loi vers
\(X\) si et seulement si
\[\lim_{n \to +\infty} \mathbb{P}(X_n = x_i ) = \mathbb{P}(X = x_i),\,\,\, \forall 1 \leq i \leq N\]
Il suffit d'écrire pour \(f\) continue bornée
\[ \mathbb{E}(f (X_n)) = \sum_{i=1}^N f (x_i ) P(X_n = x_i) \]

Dans l'exemple ci-dessus, \(N\) est fini et fixé. Mais nous avons un
résultat analogue (en faisant tendre \(N\) vers l'infini) si les
variables aléatoires ont un nombre dénombrable de valeurs. En
particulier, le cas de la convergence de la loi binomiale vers la loi de
Poisson a été traité en CPGE.
\end{exemple}

\addcontentsline{toc}{subsection}{Univers de définition}
\begin{remarque}[{Univers de
définition}]\phantomsection\label{IV-univers-de-duxe9finition}

Dans la
\protect\hyperref[IV-defconvloi]{définition (p.~\pageref*{IV-defconvloi})},
les v.a. \(X_n\) et \(X\) peuvent être définies sur des univers
distincts puisque seules leurs lois sont en cause. Il arrive même qu'une
suite \(X_n\) converge vers une limite \(X\) qui ne peut pas exister sur
les espaces sur lesquels sont définies les \(X_n\), parce que ceux-ci
sont trop ``petits''~: par exemple, si \(X_n\) est une variable
binomiale à \(n\) modalités, convenablement normalisée, et la limite
\(X\) est gaussienne (on pourra justifier ceci par le
\protect\hyperref[IV-TCL]{théorème central
limite (p.~\pageref*{IV-TCL})}); l'espace naturel sur lequel est définie
\(X_n\) contient \(n+1\) points, et sur un tel espace toutes les v.a.
sont discrètes. La convergence en loi permet donc une sorte de
convergence pour des v.a. pour lesquelles toute autre forme de
convergence serait impossible.
\end{remarque}

\runinhead{Exercice -- Suite
gaussienne (\(\mathord{\bullet}\))}\phantomsection\addcontentsline{toc}{subsubsection}{Suite
gaussienne}\label{IV-suitegauss}
Soit \((X_n)_{n \in \mathbb{N}^\ast}\) et \(X\) des variables aléatoires
de lois respectives \(\mathcal{N}(0,\sigma^2_n)\) et
\(\mathcal{N}(0,\sigma^2)\). On suppose que la suite de réels positifs
\((\sigma_n)_{n\in \mathbb{N}^\ast}\) converge vers \(\sigma > 0\) quand
\(n\) tend vers l'infini. Montrer que la suite
\((X_n)_{n \in \mathbb{N}^\ast}\) converge en loi vers
\(X\).\solref{IV-answer-suitegauss}

La convergence en loi est plus faible que la convergence en probabilité
et donc aussi que les convergences presque-sûre et en moyenne.

\addcontentsline{toc}{subsection}{En proba \(\Rightarrow\) en loi}
\begin{proposition}[{En proba \(\Rightarrow\) en
loi}]\phantomsection\label{IV-en-proba-rightarrow-en-loi}

Soient \((X_n)_{n\in\mathbb{N}\ast}\) et \(X\) des v.a., toutes définies
sur le même espace de probabilité \((\Omega, \mathcal{A}, \mathbb{P})\).
Si \(X_n \xrightarrow{\mathbb{P}} X\), alors
\(X_n\xrightarrow{\mathcal{L}} X\).
\end{proposition}

\begin{preuve}
Soit \(f\) une fonction réelle continue bornée. D'après la
\protect\hyperref[IV-propconv4]{proposition de
continuité (p.~\pageref*{IV-propconv4})}, on a
\(f(X_n) \xrightarrow{\mathbb{P}} f(X)\) et donc \(f(X_n)\) converge
aussi en moyenne vers \(f(X)\) par la
\protect\hyperref[IV-propconv2]{proposition --- cas
borné (p.~\pageref*{IV-propconv2})}. Comme
\(|\mathbb{E}(Y)|\leq \mathbb{E}(|Y|)\) pour toute variable aléatoire
réelle \(Y\), on en déduit
\(\mathbb{E}(f(X_n)) \to \mathbb{E}(f(X))\).\hfill$\blacksquare$
\end{preuve}

On peut finalement résumer les liens d'implications entre les différents
modes de convergence à l'aide de la figure ci-desous.

\begin{figure}
\centering
\includegraphics{cvgces.tex.pdf}
\caption{Relations entre modes de convergence}
\end{figure}

Un moyen efficace de caractériser la convergence en loi des variables
aléatoires réelles passe par l'étude de la suite des fonctions de
répartition.

\addcontentsline{toc}{subsection}{Convergence en loi et fonction de
répartition}
\begin{proposition}[{Convergence en loi et fonction de
répartition}]\phantomsection\label{IV-cvceloifdr}

Soient \((X_n)_{n\in\mathbb{N}^\ast}\) et \(X\) des variables aléatoires
réelles de fonctions de répartition respectives \(F_n\) et \(F\). Pour
que \(X_n \xrightarrow{\mathcal{L}} X\), il faut et il suffit que
\(F_n(x) \xrightarrow[n\to \infty]{} F(x)\) pour tout \(x\) en lequel
\(F\) est continue.

Notons que puisque la fonction \(F\) est continue à droite et
croissante, l'ensemble des points où \(F\) est continue est l'ensemble
\(D = \{x : F (x-) = F (x)\}\), et son complémentaire est au plus
dénombrable. Ainsi, \(D\) est dense\footnote{Soient \(\Omega\) un espace
  topologique et \(A\) une partie de \(\Omega\). On dit que \(A\) est
  \emph{dense} dans \(\Omega\) si l'une des propriétés équivalentes est
  satisfaite~: tout ouvert non vide de \(\Omega\) contient des éléments
  de \(A\)\,; l'adhérence de \(A\) est égale à \(\Omega\)\,; tout point
  de \(\Omega\) est adhérent à \(A\)\,; le complémentaire de \(A\) est
  d'intérieur vide.} dans \(\mathbb{R}\).
\end{proposition}

\begin{preuve}
\begin{enumerate}
\def\labelenumi{\arabic{enumi}.}
\tightlist
\item
  Supposons d'abord que \(X_n \xrightarrow{\mathcal{L}} X\). Soit
  \(a \in \mathbb{R}\) tel que \(F(a-)=F(a)\). Pour tout
  \(p \in \mathbb{N}^\ast\) et tout \(b \in \mathbb{R}\), il existe une
  fonction \(f_{p,b}\) continue bornée sur \(\mathbb{R}\) telle que
  \[ 1_{]-\infty,b]} \leq f_{p,b} \leq 1_{]-\infty,b + \frac{1}{p}]}.\]
  Alors, par définition,
  \(\mathbb{E}(f_{p,b}(X_n)) \to \mathbb{E}(f_{p,b}(X))\) quand \(n\)
  tend vers l'infini. L'inégalité ci-dessus implique que
  \(F_n(a) = \mathbb{P}(X_n \leq a) \leq \mathbb{E}(f_{p,a}(X_n))\) et
  \(\mathbb{E}(f_{p,a}(X)) \leq F(a+1/p)\)\,; donc
  \(\lim\sup_n F_n(a) \leq F(a+1/p)\) pour tout \(p\) et donc on a aussi
  \(\lim\sup_n F_n(a) \leq F(a)\). On a également que
  \(F_n(a) \geq \mathbb{E}(f_{p,a-1/p}(X_n))\) et
  \(\mathbb{E}(f_{p,a-1/p}(X)) \geq F(a-1/p)\)\,; donc
  \(\lim\inf_n F_n(a) \geq F(a-1/p)\) pour tout \(p\) et donc aussi
  \(\lim\inf_n F_n(a) \geq F(a)\), puique \(F(a-)=F(a)\). Ces deux
  résultats impliquent que \(F_n(a) \xrightarrow[n\to \infty]{} F(a)\).
\item
  Inversement, supposons \(F_n(x) \xrightarrow[n\to \infty]{} F(x)\)
  pour tout \(x\in T\), où \(T\) est une partie dense de \(\mathbb{R}\).
  Soit \(f\) une fonction continue bornée sur \(\mathbb{R}\) et
  \(\varepsilon > 0\). Soient \(a,b \in T\) avec
  \(F(a) \leq \varepsilon\) et \(F(b) \geq 1-\varepsilon\). Il existe
  \(n_0\) tel que
  \[ n\geq n_0 \Rightarrow \mathbb{P}(X_n\notin ]a,b]) = 1-F_n(b)+F_n(a) \leq 3\varepsilon.\]
  La fonction \(f\) est uniformément continue sur \([a,b]\), donc il
  existe un nombre fini de points \(a_0 = a < a_1 < \ldots < a_k = b\)
  appartenant tous à \(T\), et tels que
  \(|f(x) - f(a_i)|\leq \varepsilon\) si \(a_{i-1} \leq x \leq a_i\).
  Donc \[ g(x) = \sum_{i=1}^k f(a_i)1_{]a_{i-1},a_i]}(x)\] vérifie
  \(|f-g| \leq \varepsilon\) sur \(]a,b]\). Si \(M = \sup_x |f(x)|\), il
  vient alors
  \[|\mathbb{E}(f(X_n))-\mathbb{E}(g(X_n))| \leq M \mathbb{P}(X_n \notin [a,b]) + \varepsilon,\]
  de même pour \(X\). Enfin,
  \(\mathbb{E}(g(X_n)) = \sum_{i=1}^k f(a_i)(F_n(a_i) - F_n(a_{i-1})\),
  et de même pour \(X\) par définition de \(g\). Comme
  \((F_n(a_i))_{n \in \mathbb{N}^\ast}\) converge vers \(F(a_i)\) pour
  tout \(i\), on en déduit l'existence de \(n_1\) tel que
  \[ n\geq n_1 \Rightarrow |\mathbb{E}(g(X_n)) - \mathbb{E}(g(X))|\leq \varepsilon.\]
  Finalement, on a
  \[ n \geq \sup(n_0,n_1) \Rightarrow |\mathbb{E}(f(X_n))-\mathbb{E}(f(X))|\leq 3\varepsilon + 5 M \varepsilon .\]
  Vu l'arbitraire sur \(\varepsilon\), on en déduit que
  \(\mathbb{E}(f(X_n))\) converge vers \(\mathbb{E}(f(X))\), d'où le
  résultat.
\end{enumerate}

\hfill$\blacksquare$
\end{preuve}

\addcontentsline{toc}{subsection}{Cas où \(X\) admet une densité}
\begin{corollaire}[{Cas où \(X\) admet une
densité}]\phantomsection\label{IV-cas-ouxf9-x-admet-une-densituxe9}

Si la suite \((X_n)_{n\in\mathbb{N}^\ast}\) de variables aléatoires
réelles converge en loi vers \(X\), et si la loi de \(X\) admet une
densité, alors pour tous \(a< b\),
\[ \mathbb{P}(X_n \in ]a,b]) \xrightarrow[n\to \infty]{} \mathbb{P}(X\in]a,b]).\]
\end{corollaire}

\begin{preuve}
La fonction de répartition de \(X\) est alors continue en tout point.
(Mais pas nécessairement celle des variables aléatoires
\(X_n\).)\hfill$\blacksquare$
\end{preuve}

\section{Théorème central
limite}\label{IV-thuxe9oruxe8me-central-limite}

Ce théorème est aussi connu sous le nom de théorème de la limite
centrale. Plus simplement, il apparaît souvent sous l'abréviation TCL.

On considère une suite de variables aléatoire
\((X_n)_{n \in \mathbb{N}^\ast}\) indépendantes, de même loi et de carré
intégrable. On note \(m\) et \(\sigma^2\) l'espérance et la variance
commune aux variables \(X_n\), et \[S_n = X_1 + \ldots + X_n\] ainsi
(\(S_n = n M_n\)). On a vu que la loi des grands nombres assure que
\(M_n\) converge vers \(m\) presque-sûrement et en moyenne. On va
s'intéresser à la vitesse à laquelle cette convergence a lieu.

Pour évaluer cette vitesse, c'est-à-dire trouver un équivalent de
\(\frac{S_n}{n} - m\), on est amené à étudier la limite éventuelle de la
suite \(n^\alpha (\frac{S_n}{n} - m)\) pour différentes valeurs de
\(\alpha\)~: si \(\alpha\) est ``petit'' cette suite va encore tendre
vers 0, et elle va ``exploser'' si \(\alpha\) est ``grand''. On peut
espérer que pour une (et alors nécessairement une seule) valeur de
\(\alpha\), cette suite converge vers une limite qui n'est ni infinie ni
nulle.

Il se trouve que la réponse à cette question a un aspect ``négatif''~:
la suite \(n^\alpha (\frac{S_n}{n} - m)\) ne converge au sens
presque-sûr, ou même en probabilité, pour aucune valeur de \(\alpha\).
Elle a aussi un aspect ``positif''~: cette suite converge, au sens de la
convergence en loi, pour la même valeur \(\alpha = 1/2\) quelle que soit
la loi des \(X_n\), et toujours vers une loi normale.

Ce résultat, qui peut sembler miraculeux, a été énoncé par Laplace
(1749-1827) et démontré beaucoup plus tard par Lyapounov (1901). Il
montre le caractère universel de la loi normale en probabilités (d'où
son nom).

\addcontentsline{toc}{subsection}{central limite (TCL)}
\begin{theoreme}[{central limite (TCL)}]\phantomsection\label{IV-TCL}

Si les \(X_n\) sont des variables aléatoires réelles, indépendantes et
de même loi, de carré intégrable, d'espérance \(m\) et de variance
\(\sigma^2 >0\), alors les variables
\[ \frac{S_n -nm}{\sigma \sqrt{n}}\] convergent en loi vers une variable
aléatoire de loi \(\mathcal{N}(0,1)\).

En d'autres termes, \(\sqrt{n}(M_n - m)\) converge vers une variable
normale de loi \(\mathcal{N}(0,\sigma^2)\).
\end{theoreme}

\begin{preuve}
Soit \(\phi\) la fonction caractéristique de \(X_n - m\), et
\(Y_n = \frac{S_n -nm}{\sigma \sqrt{n}}\). Comme les \(X_n\) sont
indépendantes, la \protect\hyperref[IV-fct_carac_sum]{proposition ---
somme (p.~\pageref*{IV-fct_carac_sum})} entraîne que la fonction
caractéristique de \(Y_n\) est \begin{align*}
    \phi_{Y_n}(u) & = \phi_{\frac{1}{\sigma\sqrt{n}}\sum_{j=1}^n(X_j-m)}(u) \\
                  & = \phi_{\sum_{j=1}^n(X_j-m)}\left(\frac{u}{\sigma\sqrt{n}}\right) \\
                  & = \prod_{j=1}^n \phi \left(\frac{u}{\sigma\sqrt{n}}\right)\\
                  & = \phi\left(\frac{u}{\sigma\sqrt{n}}\right)^n
\end{align*} Comme \(\mathbb{E}(X_n -m) = 0\) et
\(\mathbb{E}((X_n-m)^2) = \sigma^2\), la
\protect\hyperref[IV-fct_carac_deriv]{proposition --- fonction
caractéristique et moments (p.~\pageref*{IV-fct_carac_deriv})} entraîne
que \[\phi'(0) = 0 \text{     et } \phi''(0) = -\sigma^2\] Si on fait le
développement de Taylor à l'ordre 2 au voisinage de zéro, on obtient
\[\phi(u) = 1 - \frac{u^2\sigma^2}{2} + u^2 h(|u|),\] où \(h(u) \to 0\)
quand \(u \to 0\). On a ainsi \begin{align*}
    \phi_{Y_n}(u) & = \phi\left(\frac{u}{\sigma\sqrt{n}}\right)^n \\
                  & = e^{n\log \phi\left(\frac{u}{\sigma\sqrt{n}}\right)} \\
                  & = e^{n\log (1 - \frac{u^2}{2n} + \frac{u^2}{n\sigma^2} h(\frac{u}{\sigma\sqrt{n}}))}
\end{align*} où \(\log\) désigne la valeur principale du logarithme
complexe\footnote{voir par exemple
  https://fr.wikipedia.org/wiki/Logarithme\_complexe} (elle vaut 0 au
point 1 et est continue dans le cercle complexe de centre 1 et de rayon
1/2 et admet le même développement limité au voisinage de \(z=1\) que le
logarithme réel). Comme \(\phi(0) = 1\) et que \(\phi\) est continue en
0, on a que pour \(u\) fixé et \(n\) assez grand,
\[\left| \phi\left(\frac{u}{\sigma\sqrt{n}}\right)-1\right| \leq 1/2.\]
En faisant \(n \to \infty\), on obtient
\[\lim_{n \to \infty} \phi_n(u) = \exp -\frac{u^2}{2}\] Le
\protect\hyperref[IV-levytheorem]{théorème de
Lévy (p.~\pageref*{IV-levytheorem})} implique alors que \(Y_n\) converge
en loi vers \(Z\) de fonction caractéristique \(\phi_Z(u) = e^{-u^2/2}\)
où l'on reconnaît la fonction caractéristique de la loi
\(\mathcal{N}(0,1)\).\hfill$\blacksquare$
\end{preuve}

\addcontentsline{toc}{subsection}{Conséquence}
\begin{remarque}[{Conséquence}]\phantomsection\label{IV-consuxe9quence}

On peut déduire de ce résultat que \(n^\alpha (\frac{S_n}{n} - m)\)
converge vers 0 (resp. \(+\infty\)) en probabilité lorsque
\(\alpha < 1/2\) (resp. \(\alpha > 1/2\)).
\end{remarque}

\subsection{Exemple~: convergence des lois
binomiales}\label{IV-exemple-convergence-des-lois-binomiales}

Supposons que \(S_n\) suive une loi binomiale \(\mathcal{B}(p,n)\). Cela
revient à dire que \(S_n\) a la même loi qu'une somme \(X_1+\ldots+X_n\)
de \(n\) variables aléatoires \(X_i\) indépendantes de loi
\(\mathcal{B}(p,1)\), i.e.~\(\mathbb{P}(X_i= 1) = p\) et
\(\mathbb{P}(X_i= 0) = 1- p\). On a alors \(m=p\) et
\(\sigma^2= p(1-p)\).

On veut calculer \(\mathbb{P}(S_n \leq x)\) pour \(x\) fixé et \(n\)
grand.

Si \(p\) est petit de sorte que \(\theta = np\) ne soit pas trop grand
(en pratique, \(\theta \leq 5\) convient), on peut utiliser
l'approximation par une loi de Poisson, vue en CPGE. Si \(p\) est très
proche de 1, de sorte que \(\theta = n(1 - p)\) soit comme ci-dessus,
alors \(n - S_n\) suit à son tour une loi proche de la loi de Poisson de
paramètre \(\theta\).

Dans les autres cas, on utilise la loi des grands nombres et le théorème
central limite~: \begin{align*}
\frac{S_n}{n} & \xrightarrow{\mathbb{P}} p,\\
\frac{S_n - np}{\sqrt{n p(1-p)}} & \xrightarrow{\mathcal{L}} Y \sim \mathcal{N}(0,1)
\end{align*}

Si on désigne par \(\Phi\) la fonction de répartition de la loi
\(\mathcal{N}(0,1)\), il vient
\[\mathbb{P}(S_n \leq x) \approx \Phi \left(\frac{x - np}{\sqrt{n p(1-p)}}\right)\]

Imaginons que l'on lance 1000 fois une pièce (non truquée). On cherche
la probabilité d'obtenir plus de 545 fois le côté Face. Le calcul exact
utilisant les lois binomiales est extrêmement lourd. Le résultat
ci-dessus nous donne une très bonne approximation. On a
\[ \mathbb{P}(S_{1000} > 545) = \mathbb{P}\left(\frac{S_{1000} - 500}{\sqrt{250}} > \frac{45}{\sqrt{250}}\right) \approx \int_{\frac{45}{\sqrt{250}}}^{+\infty} \frac{1}{\sqrt{2\pi}} e^{-x^2/2} dx.\]

Cette dernière intégrale se calcule numériquement (on trouve encore des
abaques où les valeurs de \(\Phi\) sont tabulées) et on obtient
\[ \mathbb{P}(S_{1000} > 545) \approx 1 - \Phi(2,84) \approx 0,0023.\]

\runinhead{Exercice -- Surbooking (\(\mathord{\bullet}\))}\phantomsection\addcontentsline{toc}{subsubsection}{Surbooking}\label{IV-surbooking}
400 places dans un avion, 420 billets vendus, 8\% de taux de
non-présentation~: quelle probabilité de surbooking\,? On utilisera une
approximation judicieuse.\solref{IV-answer-surbooking}

Le \protect\hyperref[IV-TCL]{théorème (p.~\pageref*{IV-TCL})} admet une
version multidimensionnelle, de preuve similaire. On considère des
vecteurs aléatoires \(X_n\) à valeurs dans \(\mathbb{R}^d\),
indépendants et de même loi, dont les composantes sont de carré
intégrable. On a un vecteur espérance \(m = E(X_n)\), et une matrice de
covariance \(C = (c_{ij} )_{i,j=1,\ldots,d}\) avec
\(c_ij = \text{Cov}(X_i,X_j)\). On peut alors énoncer le TCL
multi-dimensionnel.

\addcontentsline{toc}{subsection}{central limite multi-dimensionnel}
\begin{theoreme}[{central limite
multi-dimensionnel}]\phantomsection\label{IV-central-limite-multi-dimensionnel}

Les vecteurs aléatoires \(\frac{S_n-nm}{\sqrt{n}}\) convergent en loi
vers un vecteur aléatoire gaussien centré (i.e.~d'espérance nulle), de
matrice \(C\).
\end{theoreme}

\addcontentsline{toc}{subsection}{Vitesse de convergence}
\begin{remarque}[{Vitesse de
convergence}]\phantomsection\label{IV-vitesse-de-convergence}

Il est important de noter ici que la vitesse de convergence ne dépend
pas de la dimension des vecteurs \(X_n\).
\end{remarque}

\chapter{Annexe}\label{IV-annexe}

\section{Fonctions
caractéristiques}\label{IV-fonctions-caractuxe9ristiques}

Dans ce paragraphe, nous introduisons un outil important en calcul des
probabilités~: il s'agit de ce que l'on appelle \emph{la fonction
caractéristique} d'une variable aléatoire, et qui dans d'autres branches
des mathématiques s'appelle aussi \emph{la transformée de Fourier}. Elle
nous sera notamment nécessaire, via le théorème de Lévy, pour démontrer
le théorème central limite. L'essentiel de cette section peut être
considéré hors programme, dans le sens où, bien que très utile en
pratique, sa connaissance ne sera pas évaluée à l'examen.

On notera \(< x, y >\) le produit scalaire de deux vecteurs de
\(\mathbb{R}^n\) . Si \(u \in \mathbb{R}^n\) , la fonction (complexe)
\(x \mapsto e^{i < u,x>}\) est continue, de module 1. Donc si \(X\) est
un vecteur aléatoire à valeurs dans \(\mathbb{R}^n\) , nous pouvons
considérer \(e^{i < u,X>}\) comme une variable aléatoire à valeurs
complexes. Ses parties réelle \(Y = \cos(< u, X>)\) et imaginaire
\(Z = \sin(< u, X>)\) sont des variables aléatoires réelles. Ces
variables aléatoires réelles sont de plus bornées par 1, donc elles
admettent une espérance. Il est alors naturel d'écrire que l'espérance
de \(e^{i < u,x>}\) est
\[\mathbb{E}(e^{i < u,X>}) = \mathbb{E}(Y) + i \mathbb{E}(Z) = \mathbb{E}(\cos< u, X>) + i\mathbb{E}(\sin< u, X>) \]

\addcontentsline{toc}{subsection}{Fonction caractéristique}
\begin{definition}[{Fonction
caractéristique}]\phantomsection\label{IV-fonction-caractuxe9ristique}

Si \(X\) est un vecteur aléatoire à valeurs dans \(\mathbb{R}^n\), sa
\emph{fonction caractéristique} est la fonction \(\phi_X\) de
\(\mathbb{R}^n\) dans \(\mathbb{C}\) définie par
\[ \phi_X(u) = \mathbb{E}(e^{i < u,X>}) = \int_{\mathbb{R}^n} e^{i < u,x>} \mathbb{P}_X(dx).\]
\end{definition}

\addcontentsline{toc}{subsection}{Interprétation}
\begin{remarque}[{Interprétation}]\phantomsection\label{IV-interpruxe9tation-1}

La fonction caractéristique ne dépend en fait \textbf{que de la loi
\(\mathbb{P}_X\) de \(X\)}~: c'est la ``transformée de Fourier'' de la
loi \(\mathbb{P}_X\).

Nous verrons que cette fonction porte bien son nom, au sens où elle
caractérise la loi \(\mathbb{P}_X\). C'est une notion qui, de ce point
de vue, généralise la fonction génératrice
\(G_X(s) = \mathbb{E}(s^X), \, s \in [0,1]\), vue en CPGE dans le cas
discret. Elle vérifie \[\phi_X(u) = G_X(e^{iu}) = \mathbb{E}(e^{iuX}) \]
pour une variable \(X\) à valeurs dans \(\mathbb{N}\).
\end{remarque}

\addcontentsline{toc}{subsection}{Propriétés de la fonction
caractéristique}
\begin{proposition}[{Propriétés de la fonction
caractéristique}]\phantomsection\label{IV-propriuxe9tuxe9s-de-la-fonction-caractuxe9ristique}

\(\phi_X\) est de module inférieur à 1, continue, avec
\[ \phi(0) = 1 ;\,\,\, \phi_X(-u) = \overline{\phi_X(u)}.\]
\end{proposition}

\begin{preuve}
\(|z|\) désigne le module d'un nombre complexe \(z\).

Comme \(\mathbb{E}(Y)^2 \leq \mathbb{E}(Y^2)\) pour toute variable
réelle \(Y\), on a~:
\[ |\phi_X(u)|^2 = \mathbb{E}(\cos< u, X>)^2 + \mathbb{E}(\sin< u, X>)^2 \leq \mathbb{E}(\cos^2< u, X> + \sin^2< u, X>) = 1.\]

Pour montrer la continuité, considérons une suite
\(u_p \xrightarrow[p \to \infty]{} u\). Il y a convergence simple de
\(e^{i < u_p,X>}\) vers \(e^{i < u,X>}\). Comme ces variables aléatoires
sont de module inférieur à 1, le théorème de convergence dominée assure
que \(\phi_X(u_p) \xrightarrow[p \to \infty]{} \phi_X(u)\). \(\phi_X\)
est donc continue.\hfill$\blacksquare$
\end{preuve}

\addcontentsline{toc}{subsection}{Transformation linéaire}
\begin{proposition}[{Transformation
linéaire}]\phantomsection\label{IV-fct_carac_vec}

Si \(X\) est un vecteur aléatoire à valeurs dans \(\mathbb{R}^n\), si
\(a \in \mathbb{R}^m\) et \(A\) est une matrice réelle de taille
\(m \times n\), alors
\[ \phi_{a+AX}(u) = e^{i < u,a>} \phi_X (A^t u), \forall u \in \mathbb{R}^m\]
\end{proposition}

\begin{preuve}
Nous avons \(e^{i< u, a + AX>} = e^{i < u,a>}e^{i <A^tu, X>}\). En
effet, \(< u, A X> = < A^t u, X>\). On prend ensuite les espérances pour
obtenir le résultat.\hfill$\blacksquare$
\end{preuve}

\begin{exemple}\phantomsection\label{IV-ex}
\begin{enumerate}
\def\labelenumi{\arabic{enumi}.}
\tightlist
\item
  \(X\) suit une loi binomiale \(\mathcal{B}(n,p)\)~:
  \(\phi_X(u) = (p e^{iu}+ 1- p)^n\).
\item
  \(X\) suit une loi de Poisson \(\mathcal{P}(\theta)\)~:
  \(\phi_X(u) = e^{\theta (e^{iu}-1)}\).
\item
  \(X\) suit une loi uniforme
  \(\mathcal{U}_{[a,b]}\)~:\(\phi_X(u) = \frac{e^{iua} - e^{iub}}{iu(b-a)}\).
\item
  \(X\) suit une loi exponentielle \(\mathcal{E}(\lambda)\),
  \(\lambda > 0\)~: \(\phi_X(u) = \frac{\lambda}{\lambda - iu}\).
\item
  \(X\) suit une loi normale \(\mathcal{N}(0,1)\)~:
  \(\phi_X(u) = e^{-u^2/2}\). (calcul ci-dessous)
\item
  \(X\) suit une loi normale \(\mathcal{N}(\mu,\sigma^2)\)~:
  \(\phi_X(u) = e^{iu\mu -u^2\sigma^2/2}\). (application directe de
  \protect\hyperref[IV-fct_carac_vec]{la proposition
  ci-dessus (p.~\pageref*{IV-fct_carac_vec})})
\end{enumerate}
\end{exemple}

\subsection{Fonction caractéristique d'une variable
gaussienne}\label{IV-fctcaracgauss}

Soit \(X \sim \mathcal{N}(0,1)\). On a \begin{align*}
\phi_X(u) = \mathbb{E}(e^{iuX}) &= \int_\mathbb{R}e^{iux}\frac{1}{\sqrt{2\pi}}e^{-x^2/2}dx \\
                          &= \int_\mathbb{R}\frac{\cos(ux)}{\sqrt{2\pi}}e^{-x^2/2}dx + \int_\mathbb{R}\frac{i\sin(ux)}{\sqrt{2\pi}}e^{-x^2/2}dx
\end{align*} Comme \(x \mapsto \frac{\sin(ux)}{\sqrt{2\pi}}e^{-x^2/2}\)
est impaire et intégrable, son intégrale est nulle, d'où
\[\phi_X(u) = \int_\mathbb{R}\frac{\cos(ux)}{\sqrt{2\pi}}e^{-x^2/2}dx.\]
D'après la \protect\hyperref[IV-fct_carac_deriv]{proposition ---
fonction caractéristique et moments (p.~\pageref*{IV-fct_carac_deriv})},
on peut dériver les deux membres par rapport à \(u\), et on obtient
\[\phi'_X(u) = \frac{1}{\sqrt{2\pi}}\int_\mathbb{R}-x\sin(ux)e^{-x^2/2}dx\]
puis par intégration par parties
\[\phi'_X(u) = -\frac{1}{\sqrt{2\pi}}\int_\mathbb{R}u\cos(ux)e^{-x^2/2}dx = -u\phi_X(u).\]
Ainsi, \(\phi_X\) satisfait à l'équation différentielle
\(\phi'_X(u) = -u\phi_X(u)\), dont la solution générale est
\[\phi_X(u) = C e^{-u^2/2}.\] Comme \(\phi_X(0) = 1,\) on en déduit
finalement que \[\phi_X(u) = e^{-u^2/2}.\]

L'intérêt majeur de la fonction caractéristique réside dans le fait
qu'elle caractérise la loi de la variable aléatoire (d'où son nom).

\addcontentsline{toc}{subsection}{Caractérisation}
\begin{theoreme}[{Caractérisation}]\phantomsection\label{IV-caracfc}

La fonction caractéristique \(\phi_X\) caractérise la loi du vecteur
aléatoire \(X\). Ainsi, si deux vecteurs aléatoires \(X\) et \(Y\) ont
même fonction caractéristique, ils ont même loi.
\end{theoreme}

\begin{preuve}
Soient les fonctions suivantes avec \(\sigma > 0\)~:
\[ f_\sigma (x) = \frac{1}{(2\pi \sigma^2)^{n/2}}\exp\left(-\frac{|x|^2}{2\sigma^2}\right) \text{ et } \widehat{f}_\sigma (u)= \exp\left(-\frac{|u|^2\sigma^2}{2}\right).\]
On note que \(f_\sigma\) est la densité d'un vecteur gaussien \(Z\) de
dimension \(n\), centré et de matrice de covariance \(\sigma^2 I_n\),
autrement dit dont les composantes sont indépendantes, centrées, de
variances \(\sigma^2\).

On a \begin{align*}
\mathbb{E}(e^{i < u,Z>}) = \int f_\sigma (x) e^{i < u,x>} dx &= \int_{\mathbb{R}^n} \prod_{j=1}^n \frac{1}{\sqrt{2\pi \sigma^2}}\exp\left(-\frac{x_j^2}{2\sigma^2}+iu_j x_j\right) dx_1\ldots dx_n \\
                                  &= \prod_{j=1}^n \int_{\mathbb{R}} \frac{1}{\sqrt{2\pi \sigma^2}}\exp\left(-\frac{t^2}{2\sigma^2}+iu_j t\right) dt \\
                                  &= \prod_{j=1}^n e^{-u_j^2\sigma^2/2}\\
                                  & = \widehat{f}_\sigma(u)
\end{align*} d'après l'exemple 5 \href{ex}{ci-dessus} et en utilisant le
théorème de Fubini. On remarque ainsi que
\[ f_\sigma (u-v) = \frac{1}{(2\pi \sigma^2)^{n/2}} \widehat{f}_{\sigma}\left(\frac{u-v}{\sigma^2}\right) = \frac{1}{(2\pi \sigma^2)^{n/2}} \int f_\sigma (x) e^{i < u -v ,x>/\sigma^2} dx.\]

Supposons que \(X\) et \(Y\) admettent la même fonction caractéristique
\(\phi_X = \phi_{Y}\). En utilisant le théorème de Fubini, on a
\begin{align*}
\mathbb{E}_X(f_\sigma (X-v)) &= \int f_\sigma (u-v) \mathbb{P}_X (du) \\
                       &= \int_{\mathbb{R}^n} \frac{1}{(2\pi \sigma^2)^{n/2}} \left( \int_{\mathbb{R}^n} f_\sigma (x) e^{i < u -v ,x>/\sigma^2} dx \right) \mathbb{P}_X (du)\\
                       &= \int_{\mathbb{R}^n} f_\sigma (x) \frac{1}{(2\pi \sigma^2)^{n/2}} \left( \int_{\mathbb{R}^n}  e^{i < u ,x/\sigma^2>} \mathbb{P}_X(du) \right) e^{-i < v ,x>/\sigma^2} dx\\
                       &= \int_{\mathbb{R}^n} f_\sigma (x) \frac{1}{(2\pi \sigma^2)^{n/2}} \phi_X\left(\frac{x}{\sigma^2}\right) e^{-i < v ,x>/\sigma^2} dx,
\end{align*}

et la même égalité reste vraie pour \(\mathbb{P}_{Y}\). On en déduit que
\[\mathbb{E}(g(X)) = \int g(u) \mathbb{P}_X (du) = \int g(u) \mathbb{P}_Y (du) = \mathbb{E}(g(X'))\]
pour toute fonction \(g\) de l'espace vectoriel de fonction engendrées
par \(u \mapsto f_\sigma (u-v)\).

D'après le théorème de Stone-Weiertrass\footnote{voir par exemple
  Simmons (\protect\hyperlink{refIV-Simmons}{2003}) p.160.}, cet espace
est dense dans l'ensemble \(C_0\) des fonctions continues sur
\(\mathbb{R}^n\) et ayant une limite nulle à l'infini, pour la topologie
de la convergence uniforme (la norme est le sup sur \(\mathbb{R}^n\)).

Par suite, \(\mathbb{E}(g(X)) = \mathbb{E}(g(X'))\) pour toute fonction
\(g \in C_0\). Comme l'indicatrice de tout ouvert est limite croissante
de fonctions de \(C_0\), on en déduit que
\(\mathbb{P}_X(A) = \mathbb{E}(1_A (X))\) est égal à
\(\mathbb{P}_{X'}(A) = \mathbb{E}(1_A (X'))\) pour tout ouvert \(A\), ce
qui implique \(P_X = P_{X'}\).\hfill$\blacksquare$
\end{preuve}

La fonction caractéristique offre également un moyen commode de
caractériser l'indépendance.

\begin{corollaire}\phantomsection\label{IV-corollaire}
Soit \(X = (X_1,\ldots,X_n)\) un vecteur aléatoire à valeurs dans
\(\mathbb{R}^n\). Ses composantes \(X_i\) sont indépendantes si et
seulement si pour tous \(u_1, \ldots, u_n \in \mathbb{R}\),
\[ \phi_X(u_1,\ldots,u_n) = \prod_{j=1}^n \phi_{X_j}(u_j) \] où
\(\phi_X\) désigne la fonction caractéristique du vecteur aléatoire
\(X\), et \(\phi_{X_j}\) celle de la composante \(X_j\) pour chaque
\(j\).
\end{corollaire}

\begin{preuve}
On a \(< u,X> = \sum_{j=1}^n u_j X_j\). Si les \(X_i\) sont
indépendantes, et comme \(e^{i < u,x>} = \prod_j e^{i u_j x_j}\), nous
obtenons directement le résultat en utilisant la
\href{Probabilité\%20II.pdf\%20\#indep_fct}{proposition} (valable dans
le cas général) qui montre que l'espérance du produit de fonctions de
variables aléatoires indépendantes est égal au produit des espérances.

Supposons inversement qu'on ait
\(\phi_X (u_1,\ldots,u_n) = \prod_{j=1}^n \phi_{X_j}(u_j)\). On peut
alors construire des variables aléatoires \(X'_j\) indépendantes, telles
que \(X'_j\) et \(X_j\) aient même loi pour tout \(j\) et donc telles
que \(\phi_{X'_j} = \phi_{X_j}\). Si \(X' = (X'_1,\ldots,X'_n)\), on a
donc \(\phi_{X'} = \phi_X\). Donc \(X\) et \(X'\) ont même loi, ce qui
entraîne que pour tous boréliens \(A_j\), on ait
\[ \mathbb{P}(\bigcap_j \{X_j \in A_j\}) = \mathbb{P}(\bigcap_j \{X'_j \in A_j\}) = \prod_j \mathbb{P}(\{X'_j \in A_j\}) = \prod_j \mathbb{P}(\{X_j \in A_j\})\]
d'où l'indépendance.\hfill$\blacksquare$
\end{preuve}

\addcontentsline{toc}{subsection}{Somme}
\begin{proposition}[{Somme}]\phantomsection\label{IV-fct_carac_sum}

Si \(X\) et \(Y\) sont deux vecteurs aléatoires indépendants à valeurs
dans \(\mathbb{R}^n\), la fonction caractéristique de la somme \(X+Y\)
est donnée par \[ \phi_{X+Y} = \phi_X\phi_Y\]
\end{proposition}

\begin{preuve}
Comme \(e^{i< u, X+Y>}=e^{i< u, X>}e^{i< u, Y>}\), il suffit d'appliquer
la \href{Probabilité\%20II.pdf\%20\#indep_fct}{proposition} déjà
utilisée dans la preuve du corollaire ci-dessus.\hfill$\blacksquare$
\end{preuve}

\subsection{Exemples~:}\label{IV-exemples}

Soient \(X\) et \(Y\) deux variables aléatoires réelles indépendantes et
\(Z = X+Y\)~:

\begin{enumerate}
\def\labelenumi{\arabic{enumi}.}
\tightlist
\item
  Si \(X\) suit la loi normale \(\mathcal{N}(m,\sigma^2)\) et \(Y\) suit
  \(\mathcal{N}(m',\sigma'^2)\), alors \(Z\) suit une loi
  \(\mathcal{N}(m+m',\sigma^2+\sigma'^2)\), d'après
  l'\protect\hyperref[IV-ex]{exemple ci-dessus (p.~\pageref*{IV-ex})}
  point 6. et la \protect\hyperref[IV-fct_carac_sum]{proposition
  ci-dessus (p.~\pageref*{IV-fct_carac_sum})}.
\item
  Si \(X\) et \(Y\) suivent des lois de Poisson de paramètres \(\theta\)
  et \(\theta'\), alors \(Z\) suit une loi de Poisson de paramètre
  \(\theta + \theta'\), d'après l'\protect\hyperref[IV-ex]{exemple
  ci-dessus (p.~\pageref*{IV-ex})} point 2. et la
  \protect\hyperref[IV-fct_carac_sum]{proposition
  ci-dessus (p.~\pageref*{IV-fct_carac_sum})}.
\item
  Si \(X\) suit une loi binomiale \(\mathcal{B}(n,p)\) et \(Y\) la loi
  biomiale \(\mathcal{B}(m,p)\), alors \(Z\) suit une loi binomiale
  \(\mathcal{B}(n+m,p)\), d'après l'\protect\hyperref[IV-ex]{exemple
  ci-dessus (p.~\pageref*{IV-ex})} point 1. et la
  \protect\hyperref[IV-fct_carac_sum]{proposition
  ci-dessus (p.~\pageref*{IV-fct_carac_sum})}.
\end{enumerate}

\addcontentsline{toc}{subsection}{Fonction caractéristique et moments}
\begin{proposition}[{Fonction caractéristique et
moments}]\phantomsection\label{IV-fct_carac_deriv}

Soit \(X\) un vecteur aléatoire de \(\mathbb{R}^n\). Si la variable
\(|X|^m\) (où \(|\cdot|\) désigne la norme euclidienne) est intégrable
pour un entier \(m\), la fonction \(\phi_X\) est \(m\) fois continûment
différentiable sur \(\mathbb{R}^n\) et pour tout choix des indices
\(i_1,\ldots, i_m\),
\[ \frac{\partial^m}{\partial u_{i_1}\ldots \partial u_{i_m}} \phi_X(u) = i^m \mathbb{E}(e^{i < u,X >}X_{i_1}\ldots X_{i_m}),\]
où les \(X_j\) sont les composantes de \(X\).
\end{proposition}

\begin{preuve}
Le résultat se démontre par application itérée du théorème de dérivation
sous le signe somme.\hfill$\blacksquare$
\end{preuve}

\begin{remarque}\phantomsection\label{IV-remarque}
En prenant \(u=0\) dans la
\protect\hyperref[IV-fct_carac_sum]{proposition ---
somme (p.~\pageref*{IV-fct_carac_sum})}, la formule permet de calculer
\(\mathbb{E}(X_{i_1}\ldots X_{i_m})\) en fonction des dérivées à
l'origine de \(\phi_X\), autrement dit de calculer tous les moments du
vecteur \(X\), s'ils existent. Par exemple, si \(X\) est à valeurs
réelles et est intégrable (respectivement de carré intégrable), on a
\[\mathbb{E}(X) = i \phi'_X(0), \,\,\,(\text{resp. } \mathbb{E}(X^2) = \phi"_X(0))\]

On donne ici la définition générale d'un vecteur gaussien, qui nous sera
utile pour la démonstration du théorème central limite.
\end{remarque}

\addcontentsline{toc}{subsection}{Vecteur Gaussien (cas général)}
\begin{definition}[{Vecteur Gaussien (cas
général)}]\phantomsection\label{IV-gauss_vec}

Un vecteur aléatoire \(X=(X_1,\ldots,X_n)\) à valeurs dans
\(\mathbb{R}^n\) est dit gaussien si toute combinaison linéaire de ses
composantes, soit \(\sum_{j=1}^n a_j X_j\) où \(a_j \in \mathbb{R}\),
suit une loi normale uni-dimensionnelle (éventuellement dégénérée, par
exemple si on prend \(a_j = 0\) pour tout j).
\end{definition}

\addcontentsline{toc}{subsection}{Cas gaussien}
\begin{theoreme}[{Cas
gaussien}]\phantomsection\label{IV-fct_carac_gauss}

\(X\) est un vecteur gaussien si et seulement si sa fonction
caractéristique s'écrit \[\phi_X(u) = e^{i< u,m> - \frac{1}{2}< u,Cu>}\]
où \(m = \mathbb{E}(X) \in \mathbb{R}^n\) et \(C\) est la matrice de
covariance de \(X\) qui est donc semi-définie positive.
\end{theoreme}

\begin{preuve}
\begin{enumerate}
\def\labelenumi{\arabic{enumi}.}
\tightlist
\item
  Supposons \(\phi_X(u) = e^{i< u,m> - \frac{1}{2}< u,Cu>}\). Pour toute
  combinaison linéaire \(Y = \sum_j a_j X_j = < a, X >\), et pour tout
  \(v \in \mathbb{R}\), on a
  \[ \phi_Y(v) = \phi_X(va) = e^{iv< a ,m> - \frac{v^2}{2}< a,Ca>}\]
  donc \(Y\) suit la loi \(\mathcal{N}(< a, m>, < a, Ca>)\).
\item
  Inversement, soit \(C\) la matrice de covariance de \(X\) et \(m\) son
  vecteur moyenne. Si \(Y = < a, X > = \sum_{j=1}^n a_j X_j\) avec
  \(a\in\mathbb{R}^n\), un calcul simple montre
  \[ \mathbb{E}(Y) = < a, m >, \mathbb{V}(Y) = < a, Ca> \] Par
  hypothèse, \(Y\) suit une loi normale donc vu le point 6. de
  l'\protect\hyperref[IV-ex]{exemple plus haut (p.~\pageref*{IV-ex})},
  sa fonction caractéristique est
  \[ \phi_Y(v) = e^{iv< a ,m> - \frac{v^2}{2}< a,Ca>} \] Mais
  \(\phi_Y(1) = \phi_{< a,X>} (1) = \mathbb{E}(e^{i < a, X>}) = \phi_X(a)\),
  d'où le résultat.
\end{enumerate}

\hfill$\blacksquare$
\end{preuve}

Le théorème suivant caractérise la convergence en loi à l'aide des
fonctions caractéristiques. C'est un critère extrêmement utile dans la
pratique.

\addcontentsline{toc}{subsection}{de Lévy}
\begin{theoreme}[{de Lévy}]\phantomsection\label{IV-levytheorem}

Soit \((X_n)_{n \in \mathbb{N}^\ast}\) une suite de vecteurs aléatoires
à valeurs dans \(\mathbb{R}^d\).

\begin{enumerate}
\def\labelenumi{\arabic{enumi}.}
\item
  Si la suite \((X_n)_{n \in \mathbb{N}^\ast}\) converge en loi vers
  \(X\), alors \(\phi_{X_n}\) converge simplement vers \(\phi_X\).
\item
  Si les \(\phi_{X_n}\) convergent simplement vers une fonction
  (complexe) \(\phi\) sur \(\mathbb{R}^d\) et si cette fonction est
  \textbf{continue} en 0, alors c'est la fonction caractéristique d'une
  variable aléatoire \(X\) et \(X_n \xrightarrow{\mathcal{L}} X\).
\end{enumerate}
\end{theoreme}

\begin{preuve}
\begin{enumerate}
\def\labelenumi{\arabic{enumi}.}
\tightlist
\item
  On remarque que \(\phi_{X_n}(u) = \mathbb{E}(g_u(X_n))\) et
  \(\phi_X(u) = \mathbb{E}(g_u(X))\) où \(g_u\) est la fonction continue
  bornée \(g_u(x) = e^{i < u,x>}\). On applique alors la
  \protect\hyperref[IV-defconvloi]{définition (p.~\pageref*{IV-defconvloi})}.
\item
  Se reporter à Jacod and Protter (\protect\hyperlink{refIV-Jacod}{2003}).
\end{enumerate}

\hfill$\blacksquare$
\end{preuve}

\chapter{Exercices}\label{IV-exercices}

\section{Inégalités de
concentration}\label{IV-inuxe9galituxe9s-de-concentration}

\runinhead{Question
1}\phantomsection\addcontentsline{toc}{subsubsection}{Question
1}\label{IV-idc1}
Soit \((X_n)_n\) une suite de variables aléatoires indépendantes et
identiquement distribuées d'espérance \(\mathbb{E}(X_i) = m\) et de
variance \(\mathbb{V}(X_i) = \sigma^2 \leq 1\). Montrer que pour tout
\(\delta \in \left]0,1\right[\), avec probabilité au moins \(1-\delta\)
on a
\[\left|\frac{1}{n}\sum_{i=1}^n X_i - m \right| \leq \sqrt{\frac{1}{\delta n}}\]

On peut trouver des bornes à décroissance beaucoup plus rapide dans des
cas particuliers.\solref{IV-answer-idc1}

\runinhead{Question
2}\phantomsection\addcontentsline{toc}{subsubsection}{Question
2}\label{IV-idc2}
On suppose désormais que les \(X_i\) sont de loi
\(\mathcal{N}(m, \sigma^2 )\).

\begin{enumerate}
\def\labelenumi{\roman{enumi})}
\item
  On pose pour \(u \in \mathbb{R}: M(u) = \mathbb{E}(e^{u(X_1 -m)})\).
  Calculer \(M (u)\).
\item
  On pose \(S_n = X_1 + \ldots + X_n\) . Montrer que
  \(\forall a \in \mathbb{R}, \forall u \in \mathbb{R}_+\),
  \[\mathbb{P}(S_n - nm > a) \leq e^{-ua} (M(u))^n,\] et que
  \(\forall a \in \mathbb{R}, \forall u \in \mathbb{R}_-\),
  \[\mathbb{P}(S_n - nm < -a) \leq e^{ua} (M(u))^n.\]
\item
  Soit \(Y_n =\frac{S_n}{n}\). Montrer que \(\forall \varepsilon >0\),
  \[\mathbb{P}(|Y_n - m| > \varepsilon) \leq 2 \exp\left(\frac{-n\varepsilon^2}{2\sigma^2}\right).\]
  Cette inégalité est appelée inégalité de Chernov.
\end{enumerate}

On peut dériver ce type d'inégalités pour différentes lois de
probabilité. On voit qu'ici la concentration auprès de la moyenne se
fait à vitesse exponentielle. On a le même type de résultats pour la loi
de Bernoulli par exemple ce qui est très utile en apprentissage
statistique dans les problèmes de classification.\solref{IV-answer-idc2}

\runinhead{Question
3}\phantomsection\addcontentsline{toc}{subsubsection}{Question
3}\label{IV-idc3}

\begin{enumerate}
\def\labelenumi{\roman{enumi})}
\setcounter{enumi}{3}
\item
  On suppose que \(m = 1\) et que \(\sigma^2 = 10\). Quelle taille
  d'échantillon doit-on choisir pour obtenir
  \[\mathbb{P}(|Y_n - m| < \varepsilon) \geq \alpha,\] avec
  \(\alpha = 0,95\) et \(\varepsilon = 0,05\),

  \begin{itemize}
  \tightlist
  \item
    en utilisant l'inégalité de Bienaymé-Chebyshev\,?
  \item
    en utilisant l'inégalité de Chernov\,?
  \end{itemize}
\end{enumerate}

\solref{IV-answer-idc3}

\section{Lemme de Borel-Cantelli}\label{IV-lemme-de-borel-cantelli}

Soit \(A_n\) une suite d'événements sur l'espace probabilisé
\((\Omega,\mathcal{A}, \mathbb{P})\).

\runinhead{Question
1}\phantomsection\addcontentsline{toc}{subsubsection}{Question
1}\label{IV-bc1}
On suppose que \(\sum_{i=1}^n \mathbb{P}(A_n) < \infty\). Montrer que
\(\mathbb{P}(\lim \sup_{n \to \infty} A_n) = \mathbb{P}\left(\bigcap_{n\geq 1} \bigcup_{k \geq n} A_k \right) =0\).\solref{IV-answer-bc1}

\runinhead{Question
2}\phantomsection\addcontentsline{toc}{subsubsection}{Question
2}\label{IV-bc2}
On suppose maintenant que les événements \(A_n\) sont mutuellement
indépendants. Montrer que si \(\sum_{i=1}^n \mathbb{P}(A_n) = \infty\),
alors on a
\(\mathbb{P}(\lim \sup_{n \to \infty} A_n) = 1\).\solref{IV-answer-bc2}

\runinhead{Question
3}\phantomsection\addcontentsline{toc}{subsubsection}{Question
3}\label{IV-bc3}
Donner un exemple où \(\sum_{i=1}^n \mathbb{P}(A_n) = \infty\) et
\(\mathbb{P}(\lim \sup_{n \to \infty} A_n) < 1\) quand les \(A_n\) ne
sont pas indépendants.\solref{IV-answer-bc3}

\section{Convergence vers une
constante}\label{IV-convergence-vers-une-constante}

Soient \(\left(X_n\right)_{n\in\mathbb{N}^\ast}\) une suite de variables
aléatoires réelles, \(X\) une autre variable aléatoire réelle et
\(a \in \mathbb{R}\).

\runinhead{Convergence
\(\mathcal{L}^2\)}\phantomsection\addcontentsline{toc}{subsubsection}{Convergence
\(\mathcal{L}^2\)}\label{IV-cvtocst-l2}
On suppose que \(X\) et chaque \(X_n\), \(n\in\mathbb{N}^\ast\), sont de
carré intégrable.

\begin{enumerate}
\def\labelenumi{\arabic{enumi}.}
\tightlist
\item
  Montrer l'implication
\end{enumerate}

\[X_n \xrightarrow[n\to+\infty]{\mathcal{L}^2} X \Rightarrow \left|\begin{array}{ll} \mathbb{E}\left(X_n\right) \xrightarrow[n\to+\infty]{} \mathbb{E}(X),\\ \mathbb{V}\left(X_n\right) \xrightarrow[n\to+\infty]{} \mathbb{V}(X). \end{array}\right.\]

\begin{enumerate}
\def\labelenumi{\arabic{enumi}.}
\setcounter{enumi}{1}
\tightlist
\item
  Vérifier l'équivalence
  \[X_n \xrightarrow[n\to+\infty]{\mathcal{L}^2} a \Leftrightarrow \left|\begin{array}{ll} \mathbb{E}\left(X_n\right) \xrightarrow[n\to+\infty]{} a,\\ \mathbb{V}\left(X_n\right) \xrightarrow[n\to+\infty]{} 0. \end{array}\right.\]
\end{enumerate}

\solref{IV-answer-cvtocst-l2}

\runinhead{Convergences en loi et
probabilité}\phantomsection\addcontentsline{toc}{subsubsection}{Convergences
en loi et probabilité}\label{IV-cvtocst-loiprob}
Montrer que si \(X_n\) converge en loi vers \(a\) quand \(n\to+\infty\),
alors elle converge aussi en probabilité vers
\(a\).\solref{IV-answer-cvtocst-loiprob}

\section{Fonction de répartition
empirique}\label{IV-fonction-de-ruxe9partition-empirique}

Soient \(X\) une variable aléatoire réelle de fonction de répartition
\(F\) et \(\left(X_n\right)_{n\in\mathbb{N}^\ast}\) une suite de
variables aléatoires indépendantes, de même loi que \(X\). Pour tout
\(n\in\mathbb{N}^\ast\) on définit la fonction de répartition empirique
comme suit~:
\[F_n : x \in \mathbb{R}\mapsto \dfrac{1}{n}\,\sum_{i = 1}^n 1_{]-\infty,x]}(X_i).\]

Son nom est issu de la modélisation probabiliste en statistique. Lorsque
l'on souhaite étudier une variable physique dont on n'est pas à même de
prédire parfaitement les valeurs (e.g.~la hauteur d'eau de la Seine à la
station d'Austerlitz), le statisticien va la considérer comme une
variable aléatoire \(X\). Pour en retrouver les propriétés, il va
observer un certain nombre \(n\in\mathbb{N}^\ast\) de réalisations de
cette variable (en réalisant des mesures sur le terrain), qu'il va à
nouveau considérer comme des variables aléatoires \(X_1,\dots,X_n\) de
même loi que \(X\). Lorsque cette hypothèse est raisonnable, il va
supposer que ces dernières sont même indépendantes. Nous allons voir
dans cet exercice que pour tout \(x\in\mathbb{R}\), \(F_n(x)\), calculée
à partir des données de terrain (d'où le nom attribué à \(F_n\)),
approche bien la quantité théorique \(F(x) = \mathbb{P}(X\leq x)\).

On considère \(x\in\mathbb{R}\) fixé.

\runinhead{Préliminaire}\phantomsection\addcontentsline{toc}{subsubsection}{Préliminaire}\label{IV-fdremp-loi}
Quelle est la loi de la variable aléatoire \(1_{]-\infty,x]}(X)\)\,?
Expliciter son espérance et sa variance.\solref{IV-answer-fdremp-loi}

\runinhead{Biais}\phantomsection\addcontentsline{toc}{subsubsection}{Biais}\label{IV-fdremp-biais}
Soit \(n\in\mathbb{N}^\ast\). Calculer l'espérance de \(F_n(x)\). On dit
que \(F_n(x)\) est un \emph{estimateur sans biais} de
\(F(x)\).\solref{IV-answer-fdremp-biais}

\runinhead{Erreur quadratique
moyenne}\phantomsection\addcontentsline{toc}{subsubsection}{Erreur
quadratique moyenne}\label{IV-fdremp-mse}
Montrer que \(F_n(x) \overset{\mathcal{L}^2}{\longrightarrow} F(x)\)
quand \(n\to+\infty\).\solref{IV-answer-fdremp-mse}

\runinhead{Consistance}\phantomsection\addcontentsline{toc}{subsubsection}{Consistance}\label{IV-fdremp-ps}
Montrer que \(F_n(x) \to F(x)\) p.s. quand \(n\to+\infty\). On dit que
\(F_n(x)\) est un \emph{estimateur fortement consistant} de
\(F(x)\).\solref{IV-answer-fdremp-ps}

\section{Théorème de Weierstrass sur
\([0,1]\)}\label{IV-thuxe9oruxe8me-de-weierstrass-sur-01}

Le théorème de Weierstrass est un résultat classique d'analyse dont
l'énoncé est le suivant~: soient \(I\) un segment de \(\mathbb{R}\) et
\(f : I \to \mathbb{R}\) une fonction continue. Alors il existe une
suite \((P_n)_{n\in\mathbb{N}^\ast}\) de polynômes
\(\mathbb{R}\to\mathbb{R}\) convergeant uniformément vers \(f\) sur
\(I\).

Dans cet exercice, nous allons voir une démonstration constructive de ce
théorème dans le cas particulier où \(I = [0,1]\).

\runinhead{Preuve du
théorème.}\phantomsection\addcontentsline{toc}{subsubsection}{Preuve du
théorème.}\label{IV-weier-thm}
Soit \(\left(X_n\right)_{n\in\mathbb{N}^\ast}\) une suite de variable
aléatoires indépendantes, suivant toutes la même loi de Bernoulli de
paramètre \(x \in ]0,1[\). Pour tout \(n\in\mathbb{N}^\ast\) on pose
\(M_n := \frac{1}{n}\,\sum_{i = 1}^n X_i\). On considère une fonction
\(f : [0,1] \to \mathbb{R}\) continue.

\begin{enumerate}
\def\labelenumi{\arabic{enumi}.}
\item
  Soit \(n\in\mathbb{N}^\ast\). Vérifier que
  \(\mathbb{E}\left( f(M_n) \right)\) peut s'écrire sous forme de
  polynôme.
\item
  Montrer que
  \(\mathbb{E}\left( f(M_n) \right) \xrightarrow[n\to+\infty]{} f(x).\)
\item
  En déduire qu'il existe une suite de polynômes
  \([0,1] \to \mathbb{R}\) qui convergent simplement vers \(f\).
\item
  Montrer que la suite de polynômes exhibée à la question précédente
  converge même uniformément vers \(f\).
\end{enumerate}

\solref{IV-answer-weier-thm}

\section{Théorème de Slutsky}\label{IV-thuxe9oruxe8me-de-slutsky}

\runinhead{Question
1}\phantomsection\addcontentsline{toc}{subsubsection}{Question
1}\label{IV-slut1}
Soient \((X_n)_n\) et \(X\) des variables aléatoires. Montrer que
\(X_n \xrightarrow{\mathcal{L}} X\) si et seulement si
\(\mathbb{E}(f(X_n)) \xrightarrow[n\to\infty]{} \mathbb{E}(f(X))\) pour
toute fonction \(f\) lipschitzienne bornée.\solref{IV-answer-slut1}

\runinhead{Question
2}\phantomsection\addcontentsline{toc}{subsubsection}{Question
2}\label{IV-slut2}
Soient \((X_n)_n\) et \((Y_n)_n\) deux suites de variables aléatoires
définies sur le même espace de probabilité. On suppose que \((X_n)_n\)
converge en loi vers \(X\) et que \((X_n - Y_n)_n\) converge vers 0 en
probabilité. Montrer que \((Y_n)_n\) converge en loi vers
\(X\).\solref{IV-answer-slut2}

\runinhead{Question
3}\phantomsection\addcontentsline{toc}{subsubsection}{Question
3}\label{IV-slut3}
Montrer ce résultat à l'aide du
\protect\hyperref[IV-levytheorem]{théorème de
Lévy (p.~\pageref*{IV-levytheorem})}.

\begin{center}\rule{0.5\linewidth}{0.5pt}\end{center}

\solref{IV-answer-slut3}


\chapter{Références}\label{IV-ruxe9fuxe9rences}

\hypertarget{refs}{}
\begin{cslreferences}
\leavevmode\hypertarget{refIV-Jacod}{}%
Jacod, J., and P. Protter. 2003. \emph{L'essentiel En Théorie Des
Probabilités}. Cassini.
\url{https://hal.archives-ouvertes.fr/hal-00104956}.

\leavevmode\hypertarget{refIV-Simmons}{}%
Simmons, George F. 2003. \emph{Introduction to Topology and Modern
Analysis}. Reprint edition. Malabar, Fla: Krieger Publishing Company.
\end{cslreferences}
